\documentclass{article}
\usepackage[margin=0.8in]{geometry}
\usepackage{amsmath,amssymb,amsthm,mathtools}
\usepackage{mathrsfs,bm,bbm}
\usepackage{graphicx,booktabs,array,multirow,tabularx,longtable}
\usepackage{enumitem}
\usepackage{xcolor}
\usepackage{algorithm,algpseudocode}
\usepackage{subcaption,tikz}
\usepackage{siunitx,cancel,accents}
\usepackage{microtype}
\usepackage[numbers,sort&compress]{natbib}
\usepackage[colorlinks=true,linkcolor=blue,citecolor=blue,urlcolor=blue]{hyperref}
\usepackage{cleveref}
\setlength{\emergencystretch}{2em}
\newtheorem{assumption}{Assumption}
\newtheorem{definition}{Definition}
\newtheorem{theorem}{Theorem}
\newtheorem{lemma}{Lemma}
\newtheorem{proposition}{Proposition}
\newtheorem{openendedquestion}{Open-ended Question}
\newtheorem{conjecture}{Conjecture}
\newcommand{\coreref}[1]{\ref{#1}}

\begin{document}

\sloppy

\section*{stat\_\allowbreak{}ctc\_\allowbreak{}tailindex\_\allowbreak{}collision\_\allowbreak{}frontier}

\noindent\textit{Specialization:} v1\par

\paragraph{TL;DR.} The finite-threshold CTC contrast changes on the logarithmic collision scale \(\Delta_n\log(n/k_n)\), while tail ordering is resolved on the root-\(k_n\) scale.  Exact-Pareto and class-level forward and reverse population curves make this transition explicit without asserting that structural orientation is a unique functional of the observed law.  A local experiment indexed only by \((\lambda,h,b_1,b_2)\) is invalid because it omits root-\(k_n\) common-index variation, and its summed marginal information omits an order-\(k_n\) joint exceedance-overlap covariance.  The constant full-alphabet direction set is exactly honest and label-equivariant; the augmented paired experiment, a nonconstant envelope-conditional honest set, and a matching ambiguity frontier remain open.

\section{Project justification}

\paragraph{Gap.} Published CTC inference covers fixed common-index and fixed heterogeneous-index regimes, known-confounder-adjusted parametric estimation and permutation testing, and stationary lagged or maximum CTC.  It does not provide the joint local tail-rank experiment required for a static independent triangular-array collision with marked-mixture drift, second-order bias, common-index movement, and joint exceedance overlap, nor a matching ambiguity frontier for that experiment.

\paragraph{Niche.} The propagated-cause mixture depends on \(\Delta_n\log(n/k_n)\), whereas tail-order discrimination depends on \(\sqrt{k_n}\Delta_n\).  Their exact relation \(h_n=\kappa_n\lambda_n\) produces three constrained local faces, while the data also retain a root-\(k_n\) common-index coordinate and an order-\(k_n\) overlap intensity that the four-coordinate summed-marginal handle does not encode.

\paragraph{Fill.} The note derives the exact-Pareto and class-level forward and reverse finite-threshold collision curves, characterizes the feasible scale geometry, proves the two independent obstructions to the proposed four-coordinate summed-marginal experiment, and certifies the constant full-alphabet direction set as an exactly honest label-equivariant fallback.  It does not claim a nontrivial ambiguity lower bound or a matching frontier.

\section{Related work}

Gnecco et al. (2021, Lemma 1 and Theorems 1--2) provide fixed-common-index comparable-tail CTC theory; they do not establish the unequal-index triangular-array balance condition used here. The ratio-free exact-Pareto benchmark enlarges their equality point to a local curve, while separate class-level forward and label-exchanged reverse results establish broader sequential population maps under this note's assumptions; none is an inference theorem or an observational identification result. Leimenstoll and Schienle (2026, Proposition 3.1) give fixed heterogeneous-tail endpoints, which the exact-Pareto benchmark strictly refines locally. Their Theorem 4.1 is the independence and light-confounder CLT, Theorem 4.3 concerns fixed-fraction empirical-copula inference, and Theorem 4.4 is the gap-dependent intermediate-sequence heterogeneous-tail CLT with proof in Appendix A.6; these results inform, but do not supply, the joint collision experiment. Pasche, Chavez-Demoulin, and Davison (2023, Sections 2--3 and 5) remove much of a known confounder's effect, construct a generalized-Pareto parametric CTC estimator, and give a permutation test; their object is not an honest equality-boundary set and they do not study a static local index collision, its marked tail-rank experiment, or its minimax ambiguity loss. Bodik, Paluš, and Pawlas (2024, Definition 1, Theorems 1--2, Section 4, and Section 6) develop a lagged maximum CTC and estimator for stationary Granger-type heavy-tailed VAR and nonlinear autoregressive time series; they do not study differing-tail independent collisions, and explicitly leave distributional theory, bias correction, and rigorous testing open. Beirlant, Bouquiaux, and Werker (2006, Theorems 2.1--2.2) inform the root-\(k_n\) score and convolution argument. Drees (1998, Theorem 2.1) informs the second-order tail-process approximation. Bücher and Dette (2013, Theorem 2.2 and Section 3) inform rank-process and multiplier calibration. Carpentier and Kim (2014, Theorems 3.2 and 3.6) and Sasaki and Wang (2024) inform scalar honest adaptation and declared worst-case bias control. Novak (2014) informs the heavy-tail testing lower-bound route.

\section{Setup}

Target (estimand / structural parameter): \(D_n(t_n)=\Gamma_{12,n}(t_n)-\Gamma_{21,n}(t_n)\).
Primary functional / estimator / diagnostic: The empirical contrast is \(\widehat D_n=\widehat\Gamma_{12,n}-\widehat\Gamma_{21,n}\). For qualifying forward arrays with \(\xi_n\to\xi\), \(D_n(t_n)\to1/[2\{1+\beta_0^{\alpha_0}c_0\exp(-\xi)\}]\). For qualifying reverse arrays with \(\xi_{\mathrm{rev},n}\to\xi_{\mathrm{rev}}\), \(D_n(t_n)\to-1/[2\{1+\beta_0^{\alpha_0}c_0^{-1}\exp(-\xi_{\mathrm{rev}})\}]\). The \(\lambda\)-special cases are \(g_{\alpha_0,\beta_0,c_0}(\lambda)\) and \(-g_{\alpha_0,\beta_0,1/c_0}(-\lambda)\). On the scaled exact-Pareto subclasses, the endpoints run from \(0\) to \(1/2\) in the forward orientation and from \(-1/2\) to \(0\) in reverse as \(\lambda\) runs from \(-\infty\) to \(+\infty\). These are sequential structural-array population maps, not observed-law orientation identification and not an inference procedure.
\begin{itemize}
  \item \texttt{n} --- \texttt{positive integer}; \(\mathbb N\); \texttt{sample-\allowbreak{}size index}
  \item \texttt{k\_\allowbreak{}n} --- \texttt{positive integer sequence}; \(\{1,\ldots,n-1\}\); \texttt{intermediate threshold index}
  \par\smallskip\noindent\emph{Definition.} number of upper order statistics retained at sample size \(n\)
  \item \texttt{t\_\allowbreak{}n} --- \texttt{positive real sequence}; \texttt{tail extrapolation scale}
  \par\smallskip\noindent\emph{Definition.} \(t_n=n/k_n\)
  \item \texttt{u\_\allowbreak{}n} --- \texttt{positive real sequence}; \texttt{deterministic exceedance level for the exact-\allowbreak{}Pareto calculation}
  \item \texttt{L} --- \texttt{positive integer}; \(\mathbb N\); \texttt{fixed finite threshold-\allowbreak{}grid cardinality}
  \item \texttt{alpha\_\allowbreak{}0} --- \texttt{positive real}; \((0,\infty)\); \texttt{common limiting innovation tail index}
  \item \texttt{c\_\allowbreak{}0} --- \texttt{positive real}; \((0,\infty)\); \texttt{limiting ratio of threshold-\allowbreak{}normalized innovation tail constants}
  \item \texttt{o\_\allowbreak{}n} --- \texttt{orientation label sequence}; \(\{\mathsf{forward},\mathsf{reverse}\}\); structural orientation, with \(\mathsf{forward}\) denoting \(X_{1,n}\to X_{2,n}\)
  \item \texttt{beta\_\allowbreak{}0} --- \texttt{positive real}; \((0,\infty)\); \texttt{limiting structural coefficient}
  \item \texttt{beta\_\allowbreak{}n} --- \texttt{positive real sequence}; \((0,\infty)\); \texttt{coefficient on the oriented structural edge}
  \item \texttt{epsilon\_\allowbreak{}1n} --- \texttt{nonnegative random variable sequence}; \(\mathbb R_+\); \texttt{first innovation}
  \item \texttt{epsilon\_\allowbreak{}2n} --- \texttt{nonnegative random variable sequence}; \(\mathbb R_+\); \texttt{second innovation}
  \item \texttt{Q\_\allowbreak{}1n} --- \texttt{probability-\allowbreak{}law sequence}; \(\mathcal P(\mathbb R_+)\)
  \par\smallskip\noindent\emph{Definition.} law of \(\epsilon_{1,n}\)
  \item \texttt{Q\_\allowbreak{}2n} --- \texttt{probability-\allowbreak{}law sequence}; \(\mathcal P(\mathbb R_+)\)
  \par\smallskip\noindent\emph{Definition.} law of \(\epsilon_{2,n}\)
  \item \texttt{alpha\_\allowbreak{}1n} --- \texttt{positive real sequence}
  \par\smallskip\noindent\emph{Definition.} upper-tail index of \(Q_{1,n}\)
  \item \texttt{alpha\_\allowbreak{}2n} --- \texttt{positive real sequence}
  \par\smallskip\noindent\emph{Definition.} upper-tail index of \(Q_{2,n}\)
  \item \texttt{Delta\_\allowbreak{}n} --- \texttt{real sequence}; \texttt{signed local tail-\allowbreak{}index gap}
  \par\smallskip\noindent\emph{Definition.} \(\Delta_n=\alpha_{1,n}-\alpha_{2,n}\)
  \item \texttt{rho\_\allowbreak{}1} --- \texttt{negative real}; \(( -\infty,0)\); \texttt{first second-\allowbreak{}order exponent}
  \item \texttt{rho\_\allowbreak{}2} --- \texttt{negative real}; \(( -\infty,0)\); \texttt{second second-\allowbreak{}order exponent}
  \item \texttt{A\_\allowbreak{}1n} --- \texttt{real-\allowbreak{}valued function sequence}; \(A_{1,n}:(1,\infty)\to\mathbb R\); \texttt{first second-\allowbreak{}order auxiliary function}
  \item \texttt{A\_\allowbreak{}2n} --- \texttt{real-\allowbreak{}valued function sequence}; \(A_{2,n}:(1,\infty)\to\mathbb R\); \texttt{second second-\allowbreak{}order auxiliary function}
  \item \texttt{U\_\allowbreak{}jn} --- \texttt{indexed family of tail-\allowbreak{}quantile functions}; \(U_{j,n}:(1,\infty)\to\mathbb R_+\), \(j\in\{1,2\}\)
  \par\smallskip\noindent\emph{Definition.} \(U_{j,n}(s)=Q_{j,n}^{\leftarrow}(1-1/s)\)
  \item \texttt{c\_\allowbreak{}1n} --- \texttt{positive real sequence}; \texttt{threshold-\allowbreak{}normalized first tail constant}
  \par\smallskip\noindent\emph{Definition.} \(c_{1,n}=U_{1,n}(t_n)^{\alpha_{1,n}}/t_n\)
  \item \texttt{c\_\allowbreak{}2n} --- \texttt{positive real sequence}; \texttt{threshold-\allowbreak{}normalized second tail constant}
  \par\smallskip\noindent\emph{Definition.} \(c_{2,n}=U_{2,n}(t_n)^{\alpha_{2,n}}/t_n\)
  \item \texttt{X\_\allowbreak{}1n} --- \texttt{nonnegative random variable sequence}
  \par\smallskip\noindent\emph{Definition.} \(X_{1,n}=\epsilon_{1,n}\) under \(o_n=\mathsf{forward}\), and \(X_{1,n}=\beta_n\epsilon_{2,n}+\epsilon_{1,n}\) under \(o_n=\mathsf{reverse}\)
  \item \texttt{X\_\allowbreak{}2n} --- \texttt{nonnegative random variable sequence}
  \par\smallskip\noindent\emph{Definition.} \(X_{2,n}=\beta_n\epsilon_{1,n}+\epsilon_{2,n}\) under \(o_n=\mathsf{forward}\), and \(X_{2,n}=\epsilon_{2,n}\) under \(o_n=\mathsf{reverse}\)
  \item \texttt{P\_\allowbreak{}theta\_\allowbreak{}n} --- \texttt{probability-\allowbreak{}law sequence}; \(\mathcal P(\mathbb R_+^2)\)
  \par\smallskip\noindent\emph{Definition.} law of \((X_{1,n},X_{2,n})\) induced by the structural equations and innovation laws
  \item \texttt{theta\_\allowbreak{}n} --- \texttt{structural triangular-\allowbreak{}array tuple}; \texttt{model member containing structural and induced observational components}
  \par\smallskip\noindent\emph{Definition.} \(\theta_n=(o_n,\beta_n,Q_{1,n},Q_{2,n},P_{\theta,n},\alpha_{1,n},\alpha_{2,n},c_{1,n},c_{2,n},\rho_1,\rho_2,A_{1,n},A_{2,n},U_{1,n},U_{2,n})\)
  \item \texttt{o\_\allowbreak{}theta\_\allowbreak{}n} --- \texttt{orientation projection}; \texttt{structural orientation component of the tuple}
  \par\smallskip\noindent\emph{Definition.} \(o(\theta_n)=o_n\)
  \item \texttt{G\_\allowbreak{}n} --- \texttt{finite set of positive integers}; \texttt{reported finite threshold grid}
  \par\smallskip\noindent\emph{Definition.} \(\mathcal G_n=\{\lfloor n^{\nu_\ell}\rfloor:1\leq\ell\leq L\}\) for fixed \(L<\infty\) and \(0<\nu_1<\cdots<\nu_L<1\)
  \item \texttt{B\_\allowbreak{}n} --- \texttt{declared nonnegative function on the threshold grid}; \(B_n:\mathcal G_n\to[0,\infty)\); \texttt{conditioning envelope for second-\allowbreak{}order bias}
  \item \texttt{Theta\_\allowbreak{}CTC\_\allowbreak{}B} --- \texttt{class of structural triangular arrays}; collision model \(\Theta_{\mathrm{CTC}}(B)\) indexed by structural tuples and the declared bias envelope
  \item \texttt{theta\_\allowbreak{}par\_\allowbreak{}n} --- \texttt{structural triangular-\allowbreak{}array tuple sequence}; \texttt{forward independent exact-\allowbreak{}Pareto witness}
  \item \texttt{Z\_\allowbreak{}in} --- \texttt{random vector array}; \(\mathbb R_+^2\)
  \par\smallskip\noindent\emph{Definition.} \(Z_{i,n}=(X_{i,1,n},X_{i,2,n})\)
  \item \texttt{F\_\allowbreak{}1n} --- \texttt{distribution function sequence}; \(F_{1,n}:\mathbb R_+\to[0,1]\)
  \par\smallskip\noindent\emph{Definition.} first marginal distribution function under \(P_{\theta,n}\)
  \item \texttt{F\_\allowbreak{}2n} --- \texttt{distribution function sequence}; \(F_{2,n}:\mathbb R_+\to[0,1]\)
  \par\smallskip\noindent\emph{Definition.} second marginal distribution function under \(P_{\theta,n}\)
  \item \texttt{q\_\allowbreak{}1n} --- \texttt{positive real sequence}
  \par\smallskip\noindent\emph{Definition.} \(q_{1,n}=F_{1,n}^{\leftarrow}(1-1/t_n)\)
  \item \texttt{q\_\allowbreak{}2n} --- \texttt{positive real sequence}
  \par\smallskip\noindent\emph{Definition.} \(q_{2,n}=F_{2,n}^{\leftarrow}(1-1/t_n)\)
  \item \texttt{Gamma\_\allowbreak{}12n} --- \texttt{real sequence}
  \par\smallskip\noindent\emph{Definition.} \(\Gamma_{12,n}(t_n)=E_{P_{\theta,n}}[F_{2,n}(X_{2,n})\mid F_{1,n}(X_{1,n})>1-1/t_n]\)
  \item \texttt{Gamma\_\allowbreak{}21n} --- \texttt{real sequence}
  \par\smallskip\noindent\emph{Definition.} \(\Gamma_{21,n}(t_n)=E_{P_{\theta,n}}[F_{1,n}(X_{1,n})\mid F_{2,n}(X_{2,n})>1-1/t_n]\)
  \item \texttt{D\_\allowbreak{}n} --- \texttt{real sequence}; \texttt{finite-\allowbreak{}threshold CTC direction contrast}
  \par\smallskip\noindent\emph{Definition.} \(D_n(t_n)=\Gamma_{12,n}(t_n)-\Gamma_{21,n}(t_n)\)
  \item \texttt{hat\_\allowbreak{}D\_\allowbreak{}n} --- \texttt{real-\allowbreak{}valued statistic sequence}
  \par\smallskip\noindent\emph{Definition.} empirical-rank analogue of \(D_n(t_n)\) using the largest \(k_n\) observations in each conditioning coordinate
  \item \texttt{lambda\_\allowbreak{}n} --- \texttt{real sequence}; \texttt{mixture-\allowbreak{}collision coordinate}
  \par\smallskip\noindent\emph{Definition.} \(\lambda_n=\Delta_n\log t_n\)
  \item \texttt{h\_\allowbreak{}n} --- \texttt{real sequence}; \texttt{tail-\allowbreak{}ordering coordinate}
  \par\smallskip\noindent\emph{Definition.} \(h_n=\sqrt{k_n}\Delta_n\)
  \item \texttt{b\_\allowbreak{}1n} --- \texttt{real sequence}
  \par\smallskip\noindent\emph{Definition.} \(b_{1,n}=\sqrt{k_n}A_{1,n}(t_n)\)
  \item \texttt{b\_\allowbreak{}2n} --- \texttt{real sequence}
  \par\smallskip\noindent\emph{Definition.} \(b_{2,n}=\sqrt{k_n}A_{2,n}(t_n)\)
  \item \texttt{kappa\_\allowbreak{}n} --- \texttt{positive real sequence}; \texttt{scale-\allowbreak{}ratio coordinate}
  \par\smallskip\noindent\emph{Definition.} \(\kappa_n=\sqrt{k_n}/\log t_n\)
  \item \texttt{kappa} --- \texttt{regime index}; \([0,\infty]\); limit of \(\kappa_n\)
  \item \texttt{L\_\allowbreak{}kappa} --- \texttt{regime-\allowbreak{}specific feasible local space}; \texttt{constraint set for finite local coordinates}
  \item \texttt{R\_\allowbreak{}n} --- \texttt{positive real sequence}; \texttt{propagated-\allowbreak{}cause to idiosyncratic tail-\allowbreak{}intensity ratio}
  \par\smallskip\noindent\emph{Definition.} \(R_n=\beta_n^{\alpha_{1,n}}(c_{1,n}/c_{2,n})\exp(-\lambda_n/\alpha_0)\)
  \item \texttt{R\_\allowbreak{}lambda} --- \texttt{positive real-\allowbreak{}valued function}; \(R:\mathbb R\to(0,\infty)\); limiting propagated-to-idiosyncratic mass ratio on the forward zero-bias Pareto \(X_2\) channel
  \par\smallskip\noindent\emph{Definition.} \(R(\lambda)=\beta_0^{\alpha_0}c_0\exp(-\lambda/\alpha_0)\)
  \item \texttt{r\_\allowbreak{}n} --- \texttt{positive function sequence}; \(r_n:(0,\infty)\to(0,\infty)\)
  \par\smallskip\noindent\emph{Definition.} \(r_n(u)=\beta_n^{\alpha_{1,n}}(c_{1,n}/c_{2,n})u^{-\Delta_n}\)
  \item \texttt{g\_\allowbreak{}curve} --- \texttt{real-\allowbreak{}valued function}; \(g_{\alpha_0,\beta_0,c_0}:\mathbb R\to(0,1/2)\)
  \par\smallskip\noindent\emph{Definition.} \(g_{\alpha_0,\beta_0,c_0}(\lambda)=\{2[1+\beta_0^{\alpha_0}c_0e^{-\lambda/\alpha_0}]\}^{-1}\)
  \item \texttt{B\_\allowbreak{}mark} --- \texttt{Borel set}; Borel subsets of \((1,\infty)\times[0,1]\); \texttt{generic measurable set for a marked exceedance intensity}
  \item \texttt{N\_\allowbreak{}an} --- \texttt{indexed family of random point measures}; point measures on \((1,\infty)\times[0,1]\), \(a\in\{1,2\}\)
  \par\smallskip\noindent\emph{Definition.} For \(b\neq a\), \(N_n^{(a)}=\sum_{i=1}^n\mathbf 1\{X_{i,a,n}>q_{a,n}\}\delta_{(X_{i,a,n}/q_{a,n},F_{b,n}(X_{i,b,n}))}\), with \(q_{a,n}\) and \(F_{a,n}\) selected from the two declared marginal sequences.
  \item \texttt{nu\_\allowbreak{}an\_\allowbreak{}theta} --- \texttt{indexed family of finite measures}; \(\nu_{n,\theta}^{(a)}\) on \((1,\infty)\times[0,1]\), \(a\in\{1,2\}\); \texttt{candidate marked exceedance intensity}
  \par\smallskip\noindent\emph{Definition.} For \(b\neq a\) and Borel \(B\), \(\nu_{n,\theta}^{(a)}(B)=nP_{\theta,n}\{(X_{a,n}/q_{a,n},F_{b,n}(X_{b,n}))\in B,\ X_{a,n}>q_{a,n}\}\).
  \item \texttt{M\_\allowbreak{}an\_\allowbreak{}theta} --- \texttt{indexed family of finite measures}; \(M_{n,\theta}^{(a)}\) on \((1,\infty)\times[0,1]\), \(a\in\{1,2\}\); normalized marked intensity, whose total sample intensity remains of order \(k_n\)
  \par\smallskip\noindent\emph{Definition.} \(M_{n,\theta}^{(a)}=k_n^{-1}\nu_{n,\theta}^{(a)}\), so \(\nu_{n,\theta}^{(a)}(B)=k_nM_{n,\theta}^{(a)}(B)\).
  \item \texttt{vartheta\_\allowbreak{}0} --- \texttt{reference local parameter}; \texttt{named reference point on a feasible local face}
  \item \texttt{dot\_\allowbreak{}ell\_\allowbreak{}ar} --- \texttt{indexed family of candidate intensity scores}
  \par\smallskip\noindent\emph{Definition.} \(\dot\ell_{a,r}=d\dot M_{\vartheta_0,r}^{(a)}/dM_{\vartheta_0}^{(a)}\), the Radon--Nikodym derivative of the intensity tangent in direction \(r\), derived by Open-ended Question 1.
  \item \texttt{central\_\allowbreak{}sequence\_\allowbreak{}n} --- \texttt{candidate random linear functional}
  \par\smallskip\noindent\emph{Definition.} \(\mathbb{\Delta}^{\mathrm{cen}}_n(r)=k_n^{-1/2}\sum_{a=1}^2\int\dot\ell_{a,r}(z)\{dN_n^{(a)}(z)-k_n\,dM_{\vartheta_0}^{(a)}(z)\}\), derived by Open-ended Question 1.
  \item \texttt{I\_\allowbreak{}rs} --- \texttt{candidate bilinear information form}
  \par\smallskip\noindent\emph{Definition.} \(I_{rs}=\sum_{a=1}^2\int\dot\ell_{a,r}(z)\dot\ell_{a,s}(z)\,dM_{\vartheta_0}^{(a)}(z)\), derived by Open-ended Question 1.
  \item \texttt{T\_\allowbreak{}n} --- \texttt{sigma-\allowbreak{}field sequence}; \texttt{CTC tail-\allowbreak{}rank information}
  \par\smallskip\noindent\emph{Definition.} sigma-field generated by \(N_n^{(1)}\), \(N_n^{(2)}\), and both marginal upper-order-statistic arrays
  \item \texttt{E\_\allowbreak{}tail\_\allowbreak{}n} --- \texttt{statistical experiment sequence}; \texttt{tail-\allowbreak{}rank experiment indexed by structural tuples}
  \item \texttt{A\_\allowbreak{}dir} --- \texttt{finite set}; \(\mathcal A_{\mathrm{dir}}=\{\mathsf{forward},\mathsf{reverse}\}\); \texttt{orientation alphabet}
  \item \texttt{C\_\allowbreak{}n} --- \texttt{random set sequence}; \(\{\{\mathsf{forward}\},\{\mathsf{reverse}\},\mathcal A_{\mathrm{dir}}\}\)
  \par\smallskip\noindent\emph{Definition.} a \(\mathcal T_n\)-measurable direction confidence set
  \item \texttt{C\_\allowbreak{}star\_\allowbreak{}n} --- \texttt{random set sequence}; \texttt{currently certified honest direction-\allowbreak{}set construction}
  \par\smallskip\noindent\emph{Definition.} \(C_n^*=\mathcal A_{\mathrm{dir}}\) for every sample, the constant full-alphabet direction set certified by \(\mathrm{thm{:}constant\text{-}honest\text{-}direction\text{-}set}\).
  \item \texttt{C\_\allowbreak{}grid} --- \texttt{finite set of positive real numbers}; \texttt{reported multiplicative bias-\allowbreak{}envelope sensitivity grid}
  \par\smallskip\noindent\emph{Definition.} \(\mathcal C=\{c_1<\cdots<c_J\}\subset(0,\infty)\) for a reported fixed \(J<\infty\), with \(1\in\mathcal C\)
  \item \texttt{B\_\allowbreak{}n\_\allowbreak{}c} --- \texttt{indexed family of nonnegative functions on the threshold grid}; \texttt{bias-\allowbreak{}envelope sensitivity family}
  \par\smallskip\noindent\emph{Definition.} \(B_n^{(c)}(k)=cB_n(k)\) for \(c\in\mathcal C\) and \(k\in\mathcal G_n\)
  \item \texttt{C\_\allowbreak{}star\_\allowbreak{}n\_\allowbreak{}c} --- \texttt{indexed family of random set sequences}; \texttt{currently certified envelope-\allowbreak{}indexed honest direction-\allowbreak{}set family}
  \par\smallskip\noindent\emph{Definition.} \(C_n^{*,(c)}=\mathcal A_{\mathrm{dir}}\) for every \(c\in\mathcal C\) and every sample, the constant full-alphabet sensitivity family certified by \(\mathrm{thm{:}constant\text{-}honest\text{-}direction\text{-}set}\).
  \item \texttt{eta} --- \texttt{real number}; \((0,1/2)\); \texttt{fixed miscoverage level}
  \item \texttt{s\_\allowbreak{}CTC\_\allowbreak{}n} --- \texttt{nonnegative function on the threshold grid}; \texttt{CTC separation}
  \par\smallskip\noindent\emph{Definition.} \(s_{\mathrm{CTC},n}(k)=\sqrt{k}|D_n(n/k)|\)
  \item \texttt{s\_\allowbreak{}idx\_\allowbreak{}n} --- \texttt{nonnegative function on the threshold grid}; \texttt{tail-\allowbreak{}order separation}
  \par\smallskip\noindent\emph{Definition.} \(s_{\mathrm{idx},n}(k)=\sqrt{k}|\Delta_n|\)
  \item \texttt{H\_\allowbreak{}eta} --- \texttt{class of confidence-\allowbreak{}set sequences}; \texttt{uniformly honest tail-\allowbreak{}rank procedures}
  \item \texttt{K} --- \texttt{compact local parameter set}; \texttt{compact subset of a regime-\allowbreak{}specific feasible local space}
  \item \texttt{K\_\allowbreak{}n} --- \texttt{set of structural tuples}; triangular-array preimage of \(K\) in \(\Theta_{\mathrm{CTC}}(B)\)
  \item \texttt{v\_\allowbreak{}eta\_\allowbreak{}K} --- number in \([0,1]\); candidate only; the certified result is the tautological lower-side bracket and constant-set upper value in \(\mathrm{thm{:}four\text{-}coordinate\text{-}ambiguity\text{-}obstruction}\)
  \par\smallskip\noindent\emph{Definition.} \texttt{unestablished candidate collision-\allowbreak{}specific ambiguity frontier}
  \item \texttt{xi\_\allowbreak{}n} --- \texttt{real sequence}; \(\mathbb R\); \texttt{intrinsic forward threshold-\allowbreak{}scale collision coordinate}
  \par\smallskip\noindent\emph{Definition.} \(\xi_n=\Delta_n\log U_{2,n}(t_n)\)
  \item \texttt{xi} --- \texttt{real number}; \(\mathbb R\); \texttt{finite limit of the intrinsic collision coordinate}
  \item \texttt{R\_\allowbreak{}xi} --- \texttt{positive real}; \((0,\infty)\); limiting propagated-to-idiosyncratic mass ratio at intrinsic coordinate \(\xi\)
  \par\smallskip\noindent\emph{Definition.} \(R_\xi=\beta_0^{\alpha_0}c_0\exp(-\xi)\)
  \item \texttt{xi\_\allowbreak{}rev\_\allowbreak{}n} --- \texttt{real sequence}; \(\mathbb R\); \texttt{intrinsic reverse threshold-\allowbreak{}scale collision coordinate}
  \par\smallskip\noindent\emph{Definition.} \(\xi_{\mathrm{rev},n}=-\Delta_n\log U_{1,n}(t_n)\)
  \item \texttt{xi\_\allowbreak{}rev} --- \texttt{real number}; \(\mathbb R\); \texttt{finite limit of the intrinsic reverse collision coordinate}
  \item \texttt{R\_\allowbreak{}rev\_\allowbreak{}xi} --- \texttt{positive real}; \((0,\infty)\); \texttt{limiting propagated-\allowbreak{}to-\allowbreak{}idiosyncratic mass ratio at the intrinsic reverse coordinate}
  \par\smallskip\noindent\emph{Definition.} \(R_{\mathrm{rev},\xi}=\beta_0^{\alpha_0}c_0^{-1}\exp(-\xi_{\mathrm{rev}})\)
\end{itemize}

\section{Assumptions}

\begin{assumption}[ass:iid-triangular-sampling]\label{ass:iid-triangular-sampling}
For every \(n\) and structural tuple \(\theta_n\), the observations \(Z_{1,n},\ldots,Z_{n,n}\) are independent with common law \(P_{\theta,n}\).
\par\smallskip\noindent\emph{Standard} (Independent triangular-array sampling, \cite{vanDerVaart1998}).
\end{assumption}

\begin{assumption}[ass:linear-scm]\label{ass:linear-scm}
Under \(o_n=\mathsf{forward}\), \(X_{1,n}=\epsilon_{1,n}\) and \(X_{2,n}=\beta_nX_{1,n}+\epsilon_{2,n}\); under \(o_n=\mathsf{reverse}\), \(X_{2,n}=\epsilon_{2,n}\) and \(X_{1,n}=\beta_nX_{2,n}+\epsilon_{1,n}\).
\par\smallskip\noindent\emph{Standard} (Bivariate positive-coefficient linear structural causal model, \cite{GneccoEtAl2021}).
\end{assumption}

\begin{assumption}[ass:innovation-independence]\label{ass:innovation-independence}
For every \(n\), \(\epsilon_{1,n}\) and \(\epsilon_{2,n}\) are independent.
\par\smallskip\noindent\emph{Standard} (Independent innovations in a linear structural causal model, \cite{GneccoEtAl2021}).
\end{assumption}

\begin{assumption}[ass:coefficient-limit]\label{ass:coefficient-limit}
The structural coefficient satisfies \(\beta_n\to\beta_0>0\).
\par\smallskip\noindent\emph{Novel.} A stable positive propagation coefficient is a proposal-specific local domain that includes fixed coefficients and slowly drifting physical transmission strengths; no compact coefficient restriction is attributed to prior work.
\end{assumption}

\begin{assumption}[ass:continuous-margins]\label{ass:continuous-margins}
The distribution functions \(F_{1,n}\) and \(F_{2,n}\) are continuous.
\par\smallskip\noindent\emph{Standard} (Continuous margins for empirical copula inference, \cite{BuecherDette2013}).
\end{assumption}

\begin{assumption}[ass:uniform-second-order-variation]\label{ass:uniform-second-order-variation}
For \(j\in\{1,2\}\), either \(A_{j,n}\equiv0\) and \(U_{j,n}(sx)/U_{j,n}(s)=x^{1/\alpha_{j,n}}\), or, uniformly for \(x\) in compact subsets of \((0,\infty)\), \[\frac{U_{j,n}(sx)/U_{j,n}(s)-x^{1/\alpha_{j,n}}}{A_{j,n}(s)}\to x^{1/\alpha_0}\frac{x^{\rho_j}-1}{\rho_j}\] as \(s\to\infty\), uniformly over the triangular array.
\par\smallskip\noindent\emph{Novel.} This triangular-array uniformization is satisfied by Pareto, Burr, and Student families with tail and second-order parameters in fixed compact subsets while allowing their two tail indices to collide.
\end{assumption}

\begin{assumption}[ass:uniform-tail-normalization]\label{ass:uniform-tail-normalization}
For every fixed \(m>1\) and \(j\in\{1,2\}\), \[\sup_{s\in[m^{-1},m]}\left|\frac{Q_{j,n}\{\epsilon>U_{j,n}(t_n)s\}}{c_{j,n}\{U_{j,n}(t_n)s\}^{-\alpha_{j,n}}}-1\right|\to0.\]
\par\smallskip\noindent\emph{Novel.} The displayed equivalence fixes the triangular-array normalization on the actual threshold range. It is compatible with the second-order quantile condition and holds exactly for Pareto tails and asymptotically for standard Hall-class tails.
\end{assumption}

\begin{assumption}[ass:tail-constant-balance]\label{ass:tail-constant-balance}
The threshold-normalized tail constants satisfy \(c_{1,n}/c_{2,n}\to c_0\in(0,\infty)\).
\par\smallskip\noindent\emph{Novel.} This proposal-specific balance condition isolates collisions in which neither innovation vanishes merely by scale. Gnecco et al. (2021) supply only fixed-common-index comparable-tail precedent, not this unequal-index triangular-array condition.
\end{assumption}

\begin{assumption}[ass:local-tail-collision]\label{ass:local-tail-collision}
The indices satisfy \(\alpha_{1,n}\to\alpha_0\), \(\alpha_{2,n}\to\alpha_0\), and \(\Delta_n\to0\).
\par\smallskip\noindent\emph{Novel.} The anchor documents threshold sensitivity, slower convergence when tail-index gaps are smaller, and at least one climate comparison for which equality is not rejected; it does not establish a financial pair as a local collision. The two-sided triangular array is therefore proposal-specific sensitivity geometry rather than an empirical classification of every application.
\end{assumption}

\begin{assumption}[ass:intermediate-sequence]\label{ass:intermediate-sequence}
The threshold sequence satisfies \(k_n\to\infty\) and \(k_n/n\to0\).
\par\smallskip\noindent\emph{Standard} (Intermediate sequence in extreme-value inference, \cite{deHaanFerreira2006}).
\end{assumption}

\begin{assumption}[ass:bounded-bias-coordinates]\label{ass:bounded-bias-coordinates}
The sequence \((b_{1,n},b_{2,n})\) is bounded in \(\mathbb R^2\).
\par\smallskip\noindent\emph{Standard} (Second-order bias on the root-\(k_n\) boundary, \cite{Drees1998}).
\end{assumption}

\begin{assumption}[ass:declared-bias-envelope]\label{ass:declared-bias-envelope}
For every \(n\), \(k\in\mathcal G_n\), and \(j\in\{1,2\}\), \(\sqrt{k}|A_{j,n}(n/k)|\leq B_n(k)\).
\par\smallskip\noindent\emph{Standard} (Declared worst-case second-order bias envelope for uniform tail inference, \cite{SasakiWang2024}).
\end{assumption}

\section{Definitions}

\begin{definition}[def:collision-model-class]\label{def:collision-model-class}
\par\smallskip\noindent\emph{Defined object.} \(\Theta_{\mathrm{CTC}}(B)\)
\par\smallskip\noindent\emph{Construction.} \(\{(\theta_n)_{n\geq1}:\theta_n=(o_n,\beta_n,Q_{1,n},Q_{2,n},P_{\theta,n},\alpha_{1,n},\alpha_{2,n},c_{1,n},c_{2,n},\rho_1,\rho_2,A_{1,n},A_{2,n},U_{1,n},U_{2,n})\text{ for every }n\geq1,\text{ and the tuple sequence satisfies }\text{ass:iid-triangular-sampling},\text{ass:linear-scm},\text{ass:innovation-independence},\text{ass:coefficient-limit},\text{ass:continuous-margins},\text{ass:uniform-second-order-variation},\text{ass:uniform-tail-normalization},\text{ass:tail-constant-balance},\text{ass:local-tail-collision},\text{ass:intermediate-sequence},\text{ass:bounded-bias-coordinates},\text{ass:declared-bias-envelope}\}\)
\par\smallskip\noindent\emph{Member properties.} \texttt{ass:\allowbreak{}iid-\allowbreak{}triangular-\allowbreak{}sampling}; \texttt{ass:\allowbreak{}linear-\allowbreak{}scm}; \texttt{ass:\allowbreak{}innovation-\allowbreak{}independence}; \texttt{ass:\allowbreak{}coefficient-\allowbreak{}limit}; \texttt{ass:\allowbreak{}continuous-\allowbreak{}margins}; \texttt{ass:\allowbreak{}uniform-\allowbreak{}second-\allowbreak{}order-\allowbreak{}variation}; \texttt{ass:\allowbreak{}uniform-\allowbreak{}tail-\allowbreak{}normalization}; \texttt{ass:\allowbreak{}tail-\allowbreak{}constant-\allowbreak{}balance}; \texttt{ass:\allowbreak{}local-\allowbreak{}tail-\allowbreak{}collision}; \texttt{ass:\allowbreak{}intermediate-\allowbreak{}sequence}; \texttt{ass:\allowbreak{}bounded-\allowbreak{}bias-\allowbreak{}coordinates}; \texttt{ass:\allowbreak{}declared-\allowbreak{}bias-\allowbreak{}envelope}.
\end{definition}

\begin{definition}[def:finite-threshold-ctc]\label{def:finite-threshold-ctc}
\par\smallskip\noindent\emph{Defined object.} \texttt{finite-\allowbreak{}threshold CTC contrast}
\par\smallskip\noindent\emph{Construction.} For \(t_n=n/k_n\), set \(\Gamma_{12,n}(t_n)=E_{P_{\theta,n}}[F_{2,n}(X_{2,n})\mid F_{1,n}(X_{1,n})>1-1/t_n]\), \(\Gamma_{21,n}(t_n)=E_{P_{\theta,n}}[F_{1,n}(X_{1,n})\mid F_{2,n}(X_{2,n})>1-1/t_n]\), and \(D_n(t_n)=\Gamma_{12,n}(t_n)-\Gamma_{21,n}(t_n)\).
\par\smallskip\noindent\emph{Inputs.} \(P_{\theta,n}\); \(k_n\).
\end{definition}

\begin{definition}[def:empirical-ctc]\label{def:empirical-ctc}
\par\smallskip\noindent\emph{Defined object.} \texttt{empirical CTC contrast}
\par\smallskip\noindent\emph{Construction.} For \(a\neq b\), average the empirical rank of coordinate \(b\) over the largest \(k_n\) observations of coordinate \(a\), producing \(\widehat\Gamma_{ab,n}\), and set \(\widehat D_n=\widehat\Gamma_{12,n}-\widehat\Gamma_{21,n}\).
\par\smallskip\noindent\emph{Inputs.} \(Z_{1,n},\ldots,Z_{n,n}\); \(k_n\).
\end{definition}

\begin{definition}[def:collision-coordinates]\label{def:collision-coordinates}
\par\smallskip\noindent\emph{Defined object.} \texttt{collision coordinate vector}
\par\smallskip\noindent\emph{Construction.} \(\vartheta_n=(\lambda_n,h_n,b_{1,n},b_{2,n},\kappa_n)=(\Delta_n\log t_n,\sqrt{k_n}\Delta_n,\sqrt{k_n}A_{1,n}(t_n),\sqrt{k_n}A_{2,n}(t_n),\sqrt{k_n}/\log t_n)\).
\par\smallskip\noindent\emph{Inputs.} \(\Delta_n\); \(k_n\); \(A_{1,n}\); \(A_{2,n}\).
\end{definition}

\begin{definition}[def:feasible-local-spaces]\label{def:feasible-local-spaces}
\par\smallskip\noindent\emph{Defined object.} \(\mathcal L_\kappa\)
\par\smallskip\noindent\emph{Construction.} \[\mathcal L_0=\{(\lambda,0,b_1,b_2):\lambda,b_1,b_2\in\mathbb R\},\qquad \mathcal L_\kappa=\{(\lambda,\kappa\lambda,b_1,b_2):\lambda,b_1,b_2\in\mathbb R\}\quad(0<\kappa<\infty),\qquad \mathcal L_\infty=\{(0,h,b_1,b_2):h,b_1,b_2\in\mathbb R\}.\] A uniform claim at regime \(\kappa\) ranges over an explicitly named compact set \(K\subset\mathcal L_\kappa\).
\par\smallskip\noindent\emph{Inputs.} \(\kappa\).
\end{definition}

\begin{definition}[def:local-array-set]\label{def:local-array-set}
\par\smallskip\noindent\emph{Defined object.} \(K_n(K,\kappa;B)\)
\par\smallskip\noindent\emph{Construction.} \(K_n(K,\kappa;B)=\{\theta_n: (\theta_m)_{m\geq1}\in\Theta_{\mathrm{CTC}}(B),\ \kappa_m\to\kappa,\ (\lambda_n,h_n,b_{1,n},b_{2,n})\in K\}\) for a named compact \(K\subset\mathcal L_\kappa\).
\par\smallskip\noindent\emph{Inputs.} \(K\); \(\kappa\).
\end{definition}

\begin{definition}[def:tail-rank-experiment]\label{def:tail-rank-experiment}
\par\smallskip\noindent\emph{Defined object.} \(\mathcal E_n^{\mathrm{tail}}(\Theta_{\mathrm{CTC}}(B))\)
\par\smallskip\noindent\emph{Construction.} The experiment \(\{(\mathcal T_n,P_{\theta,n}^{\otimes n}|_{\mathcal T_n}):\theta\in\Theta_{\mathrm{CTC}}(B)\}\), where \(\mathcal T_n\) contains \(N_n^{(1)}\), \(N_n^{(2)}\), and both marginal upper-order-statistic arrays.
\par\smallskip\noindent\emph{Inputs.} \(\mathcal T_n\); \(\Theta_{\mathrm{CTC}}(B)\).
\end{definition}

\begin{definition}[def:pareto-witness]\label{def:pareto-witness}
\par\smallskip\noindent\emph{Defined object.} \(\theta_n^{\mathrm{Par}}\)
\par\smallskip\noindent\emph{Construction.} Set \(o_n=\mathsf{forward}\), take independent scaled-Pareto innovations satisfying \(Q_{j,n}\{\epsilon>x\}=c_{j,n}x^{-\alpha_{j,n}}\) for \(x\geq c_{j,n}^{1/\alpha_{j,n}}\), and induce \(P_{\theta,n}^{\mathrm{Par}}\) by \(X_{1,n}=\epsilon_{1,n}\) and \(X_{2,n}=\beta_nX_{1,n}+\epsilon_{2,n}\).
\par\smallskip\noindent\emph{Inputs.} \(\alpha_{1,n}\); \(\alpha_{2,n}\); \(c_{1,n}\); \(c_{2,n}\); \(\beta_n\).
\end{definition}

\begin{definition}[def:honest-direction-class]\label{def:honest-direction-class}
\par\smallskip\noindent\emph{Defined object.} \(\mathcal H_\eta(\Theta_{\mathrm{CTC}}(B))\)
\par\smallskip\noindent\emph{Construction.} \[\mathcal H_\eta(\Theta_{\mathrm{CTC}}(B))=\left\{(C_n)_{n\geq1}: C_n\in\{\{\mathsf{forward}\},\{\mathsf{reverse}\},\mathcal A_{\mathrm{dir}}\},\ C_n\text{ is }\mathcal T_n\text{-measurable},\ \liminf_{n\to\infty}\inf_{\theta\in\Theta_{\mathrm{CTC}}(B)}P_{\theta,n}^{\otimes n}\{o(\theta_n)\in C_n\}\geq1-\eta\right\}.\]
\par\smallskip\noindent\emph{Inputs.} \(\eta\); \(\Theta_{\mathrm{CTC}}(B)\).
\end{definition}

\begin{definition}[def:ambiguity-criterion]\label{def:ambiguity-criterion}
\par\smallskip\noindent\emph{Defined object.} \texttt{tail-\allowbreak{}rank ambiguity criterion}
\par\smallskip\noindent\emph{Construction.} For compact \(K\subset\mathcal L_\kappa\), compare the identical pair \[\liminf_{n\to\infty}\inf_{(C_n)\in\mathcal H_\eta(\Theta_{\mathrm{CTC}}(B))}\sup_{\theta_n\in K_n(K,\kappa;B)}P_{\theta,n}^{\otimes n}\{C_n=\mathcal A_{\mathrm{dir}}\}\quad\text{and}\quad\limsup_{n\to\infty}\sup_{\theta_n\in K_n(K,\kappa;B)}P_{\theta,n}^{\otimes n}\{C_n^*=\mathcal A_{\mathrm{dir}}\},\] using only \(\mathcal T_n\)-measurable procedures and the loss \(\mathbf 1\{C_n=\mathcal A_{\mathrm{dir}}\}\).
\par\smallskip\noindent\emph{Inputs.} \(K_n(K,\kappa;B)\); \(\mathcal H_\eta(\Theta_{\mathrm{CTC}}(B))\).
\end{definition}

\begin{definition}[def:joint-experiment-handle]\label{def:joint-experiment-handle}
\par\smallskip\noindent\emph{Defined object.} \texttt{normalized marked-\allowbreak{}intensity tangent handle}
\par\smallskip\noindent\emph{Construction.} For \(b\neq a\), start from \[
u_{n,\theta}^{(a)}(B)=nP_{\theta,n}\{(X_{a,n}/q_{a,n},F_{b,n}(X_{b,n}))\in B,\ X_{a,n}>q_{a,n}\}=k_nM_{n,\theta}^{(a)}(B).\] On the forward exact-Pareto zero-bias \(X_2\) channel, split \(M_{n,\theta}^{(2)}\) into propagated and idiosyncratic components with limiting total-mass ratio \(R(\lambda)=\beta_0^{\alpha_0}c_0\exp(-\lambda/\alpha_0)\), and derive both conditional rank-mark kernels. Obtain the \(X_1\) channel and reverse orientation by the structural pushforwards and label exchange. Around a named \(\vartheta_0\) on each feasible face, seek dominated intensity tangents with \(\dot\ell_{a,r}=d\dot M_{\vartheta_0,r}^{(a)}/dM_{\vartheta_0}^{(a)}\), then target \[\mathbb{\Delta}^{\mathrm{cen}}_n(r)=k_n^{-1/2}\sum_{a=1}^2\int\dot\ell_{a,r}(z)\{dN_n^{(a)}(z)-k_n\,dM_{\vartheta_0}^{(a)}(z)\},\qquad I_{rs}=\sum_{a=1}^2\int\dot\ell_{a,r}(z)\dot\ell_{a,s}(z)\,dM_{\vartheta_0}^{(a)}(z).\]
\par\smallskip\noindent\emph{Inputs.} \(\mathcal E_n^{\mathrm{tail}}(\Theta_{\mathrm{CTC}}(B))\); \(\nu_{n,\theta}^{(a)}\); \(R(\lambda)\).
\end{definition}

\begin{definition}[def:adaptive-direction-handle]\label{def:adaptive-direction-handle}
\par\smallskip\noindent\emph{Defined object.} \texttt{finite-\allowbreak{}grid envelope-\allowbreak{}conditional direction handle}
\par\smallskip\noindent\emph{Construction.} On the reported finite grid \(\mathcal G_n\), condition on the declared envelope \(B_n\), use simultaneous multiplier bands at level \(1-\eta\), subtract \(B_n(k)\), retain thresholds whose CTC and tail-index scores are mutually stable, and return exactly one of the two singletons or \(\mathcal A_{\mathrm{dir}}\). Repeat the construction for \(B_n^{(c)}=cB_n\), \(c\in\mathcal C\), with the calibration constrained so that \(c\leq c'\) implies \(C_n^{*,(c)}\subseteq C_n^{*,(c')}\).
\par\smallskip\noindent\emph{Inputs.} \(Z_{1,n},\ldots,Z_{n,n}\); \(\mathcal G_n\); \(B_n\); \(\mathcal C\).
\end{definition}

\begin{definition}[def:ambiguity-lower-bound-handle]\label{def:ambiguity-lower-bound-handle}
\par\smallskip\noindent\emph{Defined object.} \texttt{opposite-\allowbreak{}report tail-\allowbreak{}rank contiguity handle}
\par\smallskip\noindent\emph{Construction.} Use a Le Cam two-point or Assouad family that changes the structural orientation and the sign of the local tail gap while perturbing inverse tail weights and second-order auxiliaries so that the \(\mathcal T_n\)-likelihood ratios have bounded pairwise divergence.
\par\smallskip\noindent\emph{Inputs.} \(\mathcal E_n^{\mathrm{tail}}(\Theta_{\mathrm{CTC}}(B))\); \(K_n(K,\kappa;B)\).
\end{definition}

\section{Main results}

\begin{theorem}[thm:pareto-collision-curve]\label{thm:pareto-collision-curve}
For the forward scaled exact-Pareto tuple \(\theta_n^{\mathrm{Par}}\), let \(u_n>0\) satisfy
\[
\frac{u_n}{\beta_n c_{1,n}^{1/\alpha_{1,n}}+c_{2,n}^{1/\alpha_{2,n}}}\to\infty,
\]
and let \(r_n(u)=\beta_n^{\alpha_{1,n}}(c_{1,n}/c_{2,n})u^{-\Delta_n}\). Then
\[
E_{P_{\theta,n}^{\mathrm{Par}}}[F_{1,n}(X_{1,n})\mid X_{2,n}>u_n]
=\frac12+\frac12\frac{r_n(u_n)}{1+r_n(u_n)}+o(1).
\]
If, in addition, \(\Delta_n\log c_{2,n}\to0\), then \(\Gamma_{12,n}(t_n)\to1\). If moreover \(\lambda_n\to\lambda\in[-\infty,+\infty]\), define \(R(\lambda)=\beta_0^{\alpha_0}c_0e^{-\lambda/\alpha_0}\) when \(\lambda\in\mathbb R\). Then
\[
\bigl(\Gamma_{12,n}(t_n),\Gamma_{21,n}(t_n),D_n(t_n)\bigr)
\longrightarrow
\begin{cases}
(1,1,0), & \lambda=-\infty,\\
\displaystyle\left(1,\frac12+\frac12\frac{R(\lambda)}{1+R(\lambda)},
g_{\alpha_0,\beta_0,c_0}(\lambda)\right), & \lambda\in\mathbb R,\\
\displaystyle\left(1,\frac12,\frac12\right), & \lambda=+\infty,
\end{cases}
\]
where, for finite \(\lambda\),
\[
g_{\alpha_0,\beta_0,c_0}(\lambda)
=\frac{1}{2\{1+R(\lambda)\}}
=\{2[1+\beta_0^{\alpha_0}c_0e^{-\lambda/\alpha_0}]\}^{-1}.
\]
In particular, the finite expression equals \(\{2(1+\beta_0^{\alpha_0}c_0)\}^{-1}\) at \(\lambda=0\).
\end{theorem}
\begin{proof}
Put
\[
s_{j,n}=c_{j,n}^{1/\alpha_{j,n}}\quad(j\in\{1,2\}),
\qquad A_n=\beta_n\epsilon_{1,n},
\qquad B'_n=\epsilon_{2,n}.
\tag{1}
\]
By \(\mathrm{def{:}pareto\text{-}witness}\), \(s_{j,n}\) is the lower support endpoint of \(Q_{j,n}\). Assumption \(\mathrm{ass{:}local\text{-}tail\text{-}collision}\) and \(\alpha_0>0\) give constants \(0<a_-<a_+<\infty\) such that
\[
a_-\leq\alpha_{j,n}\leq a_+
\qquad(j\in\{1,2\})
\tag{2}
\]
for all sufficiently large \(n\). Assumption \(\mathrm{ass{:}coefficient\text{-}limit}\) likewise bounds \(\beta_n\) above and away from zero.

Let \(v_n>0\) be any deterministic level sequence satisfying
\[
\frac{v_n}{\beta_ns_{1,n}+s_{2,n}}\to\infty.
\tag{3}
\]
For every fixed \(\delta\in(0,1/2)\), condition (3) implies that \(v_n\), \(\delta v_n\), and \((1-\delta)v_n\) exceed the relevant component support endpoints for all sufficiently large \(n\). The exact-Pareto tails in \(\mathrm{def{:}pareto\text{-}witness}\) therefore give
\[
p_{1,n}(v_n):=P(A_n>v_n)
=c_{1,n}\beta_n^{\alpha_{1,n}}v_n^{-\alpha_{1,n}},
\qquad
p_{2,n}(v_n):=P(B'_n>v_n)
=c_{2,n}v_n^{-\alpha_{2,n}}.
\tag{4}
\]
Indeed, \(\beta_ns_{1,n}/v_n\to0\) and \(s_{2,n}/v_n\to0\) by (3). The lower bound in (2) consequently yields
\[
p_{1,n}(v_n)=\left(\frac{\beta_ns_{1,n}}{v_n}\right)^{\alpha_{1,n}}\to0,
\qquad
p_{2,n}(v_n)=\left(\frac{s_{2,n}}{v_n}\right)^{\alpha_{2,n}}\to0.
\tag{5}
\]
This conclusion uses only the relative-support condition (3), not absolute divergence of \(v_n\).

We first establish a ratio-free one-large-jump calculation. Fix \(\delta\in(0,1/2)\), abbreviate \(p_{j,n}=p_{j,n}(v_n)\), and note the deterministic inclusions
\[
\{A_n>v_n\}\cup\{B'_n>v_n\}
\subseteq\{A_n+B'_n>v_n\}
\tag{6}
\]
and
\[
\{A_n+B'_n>v_n\}
\subseteq
\{A_n>(1-\delta)v_n\}\cup\{B'_n>(1-\delta)v_n\}
\cup\bigl(\{A_n>\delta v_n\}\cap\{B'_n>\delta v_n\}\bigr).
\tag{7}
\]
Indeed, if the left event in (7) occurs and neither of its first two right-hand events occurs, then neither summand can be at most \(\delta v_n\), because that would force the other to exceed \((1-\delta)v_n\). By \(\mathrm{ass{:}innovation\text{-}independence}\), (4), (6), and (7) imply
\[
\begin{split}
p_{1,n}+p_{2,n}-p_{1,n}p_{2,n}
&\leq P(A_n+B'_n>v_n)\\
&\leq p_{1,n}(1-\delta)^{-\alpha_{1,n}}
+p_{2,n}(1-\delta)^{-\alpha_{2,n}}
+p_{1,n}p_{2,n}\delta^{-\alpha_{1,n}-\alpha_{2,n}}.
\end{split}
\tag{8}
\]
Moreover, (5) gives
\[
0\leq\frac{p_{1,n}p_{2,n}}{p_{1,n}+p_{2,n}}
\leq\min(p_{1,n},p_{2,n})\to0.
\tag{9}
\]
Divide (8) by \(p_{1,n}+p_{2,n}\). For fixed \(\delta\), (2) and (9) bound the upper limit by a quantity tending to one as \(\delta\downarrow0\), whereas the lower limit is at least one by (9). Thus
\[
P(A_n+B'_n>v_n)
=p_{1,n}(v_n)+p_{2,n}(v_n)
+o\bigl(p_{1,n}(v_n)+p_{2,n}(v_n)\bigr).
\tag{10}
\]
This derivation is uniform over the ratio \(p_{1,n}(v_n)/p_{2,n}(v_n)\).

Let \(H_n=F_{1,n}(\epsilon_{1,n})\). Under the forward construction in \(\mathrm{def{:}pareto\text{-}witness}\), \(F_{1,n}\) is the exact-Pareto distribution function. Directly from that distribution, \(P(H_n\leq z)=z\) for \(0\leq z\leq1\), so \(0\leq H_n\leq1\) and \(E(H_n)=1/2\). On \(\{A_n>v_n\}\), (4) gives
\[
0\leq1-H_n
=c_{1,n}\epsilon_{1,n}^{-\alpha_{1,n}}
\leq c_{1,n}(v_n/\beta_n)^{-\alpha_{1,n}}
=p_{1,n}(v_n).
\tag{11}
\]
Consequently,
\[
E[H_n\mathbf 1\{A_n>v_n\}]
=p_{1,n}(v_n)+O\bigl(p_{1,n}(v_n)^2\bigr),
\qquad
E[H_n\mathbf 1\{B'_n>v_n\}]
=\frac12p_{2,n}(v_n),
\tag{12}
\]
where the second equality follows from \(\mathrm{ass{:}innovation\text{-}independence}\). The intersection of the two component exceedance events has probability \(p_{1,n}(v_n)p_{2,n}(v_n)\). By subtracting the probability of their union from (10), the part of \(\{A_n+B'_n>v_n\}\) outside that union has probability \(o(p_{1,n}(v_n)+p_{2,n}(v_n))\), using (9). Since \(0\leq H_n\leq1\), (10)--(12) yield
\[
E[H_n\mathbf 1\{A_n+B'_n>v_n\}]
=p_{1,n}(v_n)+\frac12p_{2,n}(v_n)
+o\bigl(p_{1,n}(v_n)+p_{2,n}(v_n)\bigr).
\tag{13}
\]
Under \(\mathrm{ass{:}linear\text{-}scm}\) and the forward orientation in \(\mathrm{def{:}pareto\text{-}witness}\), \(X_{1,n}=\epsilon_{1,n}\) and \(X_{2,n}=A_n+B'_n\). Dividing (13) by (10), and using
\[
\frac{p_{1,n}(v_n)}{p_{2,n}(v_n)}
=\beta_n^{\alpha_{1,n}}\frac{c_{1,n}}{c_{2,n}}v_n^{-\Delta_n}
=r_n(v_n),
\tag{14}
\]
proves
\[
E[F_{1,n}(X_{1,n})\mid X_{2,n}>v_n]
=\frac12+\frac12\frac{r_n(v_n)}{1+r_n(v_n)}+o(1).
\tag{15}
\]
The theorem's level \(v_n=u_n\) is positive and satisfies (3), so (15) proves the arbitrary-level assertion without assuming \(u_n\to\infty\) and without any condition on \(|\Delta_n|\log u_n\).

Now assume \(\Delta_n\log c_{2,n}\to0\). Assumption \(\mathrm{ass{:}intermediate\text{-}sequence}\) gives \(t_n=n/k_n\to\infty\), and the exact \((1-1/t_n)\)-quantile of \(X_{1,n}\) is
\[
q_{1,n}=s_{1,n}t_n^{1/\alpha_{1,n}}.
\tag{16}
\]
Assumption \(\mathrm{ass{:}tail\text{-}constant\text{-}balance}\) gives \(\log(c_{1,n}/c_{2,n})=O(1)\). Hence
\[
\log\frac{q_{1,n}}{s_{2,n}}
=\frac{\log t_n}{\alpha_{1,n}}
+\frac{\log(c_{1,n}/c_{2,n})}{\alpha_{1,n}}
-\frac{\Delta_n\log c_{2,n}}{\alpha_{1,n}\alpha_{2,n}}
\longrightarrow\infty.
\tag{17}
\]
Together with \(q_{1,n}/s_{1,n}=t_n^{1/\alpha_{1,n}}\to\infty\) and the bounds on \(\beta_n\), this implies
\[
\frac{\beta_nq_{1,n}}{\beta_ns_{1,n}+s_{2,n}}\to\infty.
\tag{18}
\]
Apply the ratio-free identity (10), which only requires (3), at \(v_n=\beta_nq_{1,n}\). The two component exceedance probabilities are
\[
P(A_n>\beta_nq_{1,n})=\frac1{t_n},
\qquad
P(B'_n>\beta_nq_{1,n})
=\left(\frac{s_{2,n}}{\beta_nq_{1,n}}\right)^{\alpha_{2,n}},
\tag{19}
\]
and both tend to zero by (17)--(18). Therefore
\[
P(X_{2,n}>\beta_nq_{1,n})\to0.
\tag{20}
\]
The exact-Pareto first margin is continuous, so the event \(F_{1,n}(X_{1,n})>1-1/t_n\) equals \(\{X_{1,n}>q_{1,n}\}\) up to a null set. On that event, positivity of \(\beta_n\) and nonnegativity of the second innovation give \(X_{2,n}=\beta_nX_{1,n}+\epsilon_{2,n}>\beta_nq_{1,n}\). Monotonicity of \(F_{2,n}\), (20), and \(\mathrm{def{:}finite\text{-}threshold\text{-}ctc}\) consequently give
\[
0\leq1-\Gamma_{12,n}(t_n)
\leq1-F_{2,n}(\beta_nq_{1,n})
=P(X_{2,n}>\beta_nq_{1,n})\to0.
\tag{21}
\]

The variables \(A_n\) and \(B'_n\) have positive continuous densities on the interiors of their support intervals. Their independent sum therefore has a continuous distribution function that is strictly increasing above \(\beta_ns_{1,n}+s_{2,n}\). For all sufficiently large \(n\), let \(q_{2,n}\) be its unique \((1-1/t_n)\)-quantile, so
\[
P(X_{2,n}>q_{2,n})=\frac1{t_n}.
\tag{22}
\]
Set
\[
a_n=\beta_ns_{1,n}t_n^{1/\alpha_{1,n}},
\qquad b_n=s_{2,n}t_n^{1/\alpha_{2,n}},
\qquad M_n=\max(a_n,b_n).
\tag{23}
\]
Because \(X_{2,n}\geq A_n\vee B'_n\), (22) implies \(q_{2,n}\geq M_n\). Conversely, if each component is at most its \((1-1/(2t_n))\)-quantile, their sum is at most
\[
\beta_ns_{1,n}(2t_n)^{1/\alpha_{1,n}}
+s_{2,n}(2t_n)^{1/\alpha_{2,n}}.
\]
The union bound and (22) therefore imply, using (2), that for a fixed \(C<\infty\),
\[
M_n\leq q_{2,n}
\leq\beta_ns_{1,n}(2t_n)^{1/\alpha_{1,n}}
+s_{2,n}(2t_n)^{1/\alpha_{2,n}}
\leq CM_n.
\tag{24}
\]
The lower bound in (24) also gives
\[
\frac{q_{2,n}}{\beta_ns_{1,n}+s_{2,n}}
\geq\frac12\min\{t_n^{1/\alpha_{1,n}},t_n^{1/\alpha_{2,n}}\}
\longrightarrow\infty.
\tag{25}
\]

From \(\mathrm{ass{:}tail\text{-}constant\text{-}balance}\), \(\mathrm{ass{:}local\text{-}tail\text{-}collision}\), and the imposed scale condition,
\[
\Delta_n\log c_{1,n}
=\Delta_n\log c_{2,n}
+\Delta_n\log(c_{1,n}/c_{2,n})\to0.
\tag{26}
\]
Since \(\log\beta_n\) is bounded, (23), (26), and \(\lambda_n=\Delta_n\log t_n\) yield
\[
\Delta_n\log a_n=\frac{\lambda_n}{\alpha_{1,n}}+o(1),
\qquad
\Delta_n\log b_n=\frac{\lambda_n}{\alpha_{2,n}}+o(1).
\tag{27}
\]
If \(\lambda_n\to\lambda\in[-\infty,+\infty]\), both quantities in (27) converge to \(\lambda/\alpha_0\) in the finite case and to the corresponding signed infinity at the two endpoints. If \(\Delta_n\geq0\), then \(\Delta_n\log M_n\) is the maximum of the two quantities in (27); if \(\Delta_n<0\), it is their minimum. It therefore has the same extended-real limit in either case. By (24) and \(\Delta_n\to0\),
\[
|\Delta_n|\,|\log q_{2,n}-\log M_n|
\leq|\Delta_n|\log C\to0.
\tag{28}
\]
Thus \(\Delta_n\log q_{2,n}\) has that same limit, and \(\mathrm{ass{:}coefficient\text{-}limit}\) together with \(\mathrm{ass{:}tail\text{-}constant\text{-}balance}\) gives
\[
r_n(q_{2,n})\longrightarrow
\begin{cases}
+\infty, & \lambda=-\infty,\\
\beta_0^{\alpha_0}c_0e^{-\lambda/\alpha_0}=R(\lambda), & \lambda\in\mathbb R,\\
0, & \lambda=+\infty.
\end{cases}
\tag{29}
\]
Condition (25) permits application of the ratio-free calculation (15) at \(v_n=q_{2,n}\); as noted after (10), this calculation requires only the relative-support condition and does not require \(q_{2,n}\to\infty\). Continuity and (22) identify the conditioning event in \(\Gamma_{21,n}(t_n)\) with \(\{X_{2,n}>q_{2,n}\}\) up to a null set. Substituting (29) into (15) proves
\[
\Gamma_{21,n}(t_n)\longrightarrow
\begin{cases}
1, & \lambda=-\infty,\\
\displaystyle\frac12+\frac12\frac{R(\lambda)}{1+R(\lambda)}, & \lambda\in\mathbb R,\\
\frac12, & \lambda=+\infty.
\end{cases}
\tag{30}
\]
Equation (21) proves \(\Gamma_{12,n}(t_n)\to1\) in all three cases. Subtracting (30) from (21), as prescribed by \(D_n(t_n)=\Gamma_{12,n}(t_n)-\Gamma_{21,n}(t_n)\) in \(\mathrm{def{:}finite\text{-}threshold\text{-}ctc}\), gives the asserted joint endpoint triples and, for finite \(\lambda\),
\[
D_n(t_n)\to\frac1{2\{1+R(\lambda)\}}
=\{2[1+\beta_0^{\alpha_0}c_0e^{-\lambda/\alpha_0}]\}^{-1}
=g_{\alpha_0,\beta_0,c_0}(\lambda).
\tag{31}
\]
Setting \(\lambda=0\) in (31) gives \(\{2(1+\beta_0^{\alpha_0}c_0)\}^{-1}\), completing the proof.
\end{proof}
\par\smallskip\noindent \textit{Justification.} The ratio-free exact-Pareto calculation removes an unnecessary absolute-divergence premise and establishes the actual sequential endpoint triples at \(\lambda=\pm\infty\). \textit{Gap.} Fixed-common-index and fixed-ordering results do not give this two-sided sequential exact-Pareto population transition. \textit{Consumer.} This theorem is the exact-Pareto population benchmark for finite-threshold CTC. Separate class-level forward and reverse-collision theorems establish the broader forward and reverse population maps. This theorem supplies neither inference nor observational-law identification of orientation.

\begin{proposition}[prop:pareto-witness-membership]\label{prop:pareto-witness-membership}
If its parameter sequences obey \(\beta_n\to\beta_0>0\), \(\alpha_{j,n}\to\alpha_0\), and \(c_{1,n}/c_{2,n}\to c_0\in(0,\infty)\), then the structural tuple sequence \((\theta_n^{\mathrm{Par}})_{n\geq1}\) belongs to \(\Theta_{\mathrm{CTC}}(B)\), with \(A_{1,n}=A_{2,n}=0\), for every intermediate sequence \(k_n\) and every declared nonnegative envelope \(B_n\).
\end{proposition}
\begin{proof}
Fix parameter sequences satisfying the hypotheses of the proposition, an intermediate sequence \(k_n\), and a declared envelope \(B_n:\mathcal G_n\to[0,\infty)\).  We verify each member property in \(\mathrm{def:collision\text{-}model\text{-}class}\).

For \(j\in\{1,2\}\), put \(a_{j,n}=c_{j,n}^{1/\alpha_{j,n}}>0\).  The scaled-Pareto law specified by \(\mathrm{def:pareto\text{-}witness}\) has distribution function
\[
 F_{\epsilon_j,n}(x)=
 \begin{cases}
 0,&0\leq x<a_{j,n},\\
 1-c_{j,n}x^{-\alpha_{j,n}},&x\geq a_{j,n}.
 \end{cases}
\]
The branches agree at \(a_{j,n}\), since \(c_{j,n}a_{j,n}^{-\alpha_{j,n}}=1\), so both innovation laws are continuous and atomless.  Solving \(F_{\epsilon_j,n}(x)=1-1/s\) gives, for every \(s>1\),
\[
 U_{j,n}(s)=(c_{j,n}s)^{1/\alpha_{j,n}}. \tag{1}
\]
Because \(k_n\in\{1,\ldots,n-1\}\), one has \(t_n=n/k_n>1\), and (1) yields
\[
 \frac{U_{j,n}(t_n)^{\alpha_{j,n}}}{t_n}
 =\frac{c_{j,n}t_n}{t_n}=c_{j,n}. \tag{2}
\]
Thus the scale constants used in the Pareto laws coincide with the threshold-normalized constants in the structural tuple.

By \(\mathrm{def:pareto\text{-}witness}\), \(o_n=\mathsf{forward}\), the innovations are independent, and
\[
 X_{1,n}=\epsilon_{1,n},\qquad
 X_{2,n}=\beta_nX_{1,n}+\epsilon_{2,n}. \tag{3}
\]
Equation (3) is the forward clause of \(\mathrm{ass:linear\text{-}scm}\), while the stipulated independence is \(\mathrm{ass:innovation\text{-}independence}\).  The assumed convergence \(\beta_n\to\beta_0>0\) verifies \(\mathrm{ass:coefficient\text{-}limit}\).  The first margin is continuous because \(\epsilon_{1,n}\) is atomless.  For every \(x\geq0\), independence and atomlessness of \(\epsilon_{2,n}\) imply
\[
 \Pr\!\left(\beta_n\epsilon_{1,n}+\epsilon_{2,n}=x\right)
 =E\!\left[
   \Pr\!\left(\epsilon_{2,n}=x-\beta_n\epsilon_{1,n}
      \mid\epsilon_{1,n}\right)
 \right]=0. \tag{4}
\]
Hence the second margin is also atomless, so both marginal distribution functions are continuous, as required by \(\mathrm{ass:continuous\text{-}margins}\).  For each \(n\), realize \(Z_{1,n},\ldots,Z_{n,n}\) on the product probability space with law \(\bigl(P_{\theta,n}^{\mathrm{Par}}\bigr)^{\otimes n}\).  The coordinate projections are then independent and have common law \(P_{\theta,n}^{\mathrm{Par}}\), which verifies \(\mathrm{ass:iid\text{-}triangular\text{-}sampling}\).

Set \(A_{1,n}=A_{2,n}=0\).  From (1), for every \(s>1\) and every \(x>0\) such that \(sx>1\),
\[
 \frac{U_{j,n}(sx)}{U_{j,n}(s)}=x^{1/\alpha_{j,n}}. \tag{5}
\]
If \(K\subset(0,\infty)\) is compact, then \(\inf K>0\), so \(sx>1\) for every \(x\in K\) once \(s>1/\inf K\).  Thus (5) holds eventually and uniformly on each such \(K\).  This is exactly the zero-auxiliary branch of \(\mathrm{ass:uniform\text{-}second\text{-}order\text{-}variation}\), with identically zero remainder.

Fix \(m>1\).  Since \(\alpha_{j,n}\to\alpha_0>0\), eventually \(0<\alpha_{j,n}\leq2\alpha_0\).  By \(\mathrm{ass:intermediate\text{-}sequence}\), \(t_n=n/k_n\to\infty\).  Therefore
\[
 \inf_{s\in[m^{-1},m]}
 \frac{U_{j,n}(t_n)s}{a_{j,n}}
 =m^{-1}t_n^{1/\alpha_{j,n}}
 \geq m^{-1}t_n^{1/(2\alpha_0)}\longrightarrow\infty. \tag{6}
\]
Consequently, for all sufficiently large \(n\), the Pareto survival formula applies to \(U_{j,n}(t_n)s\) simultaneously for all \(s\in[m^{-1},m]\).  Equation (2) then gives
\[
 \sup_{s\in[m^{-1},m]}
 \left|
 \frac{Q_{j,n}\{\epsilon>U_{j,n}(t_n)s\}}
 {c_{j,n}\{U_{j,n}(t_n)s\}^{-\alpha_{j,n}}}-1
 \right|=0. \tag{7}
\]
This verifies \(\mathrm{ass:uniform\text{-}tail\text{-}normalization}\).  The assumed convergence \(c_{1,n}/c_{2,n}\to c_0\in(0,\infty)\) is precisely \(\mathrm{ass:tail\text{-}constant\text{-}balance}\).

The limits \(\alpha_{1,n}\to\alpha_0\) and \(\alpha_{2,n}\to\alpha_0\) imply
\[
 \Delta_n=\alpha_{1,n}-\alpha_{2,n}\longrightarrow0, \tag{8}
\]
so \(\mathrm{ass:local\text{-}tail\text{-}collision}\) holds.  The chosen \(k_n\) obeys \(\mathrm{ass:intermediate\text{-}sequence}\) by hypothesis.  Finally, the zero auxiliary functions imply
\[
 b_{j,n}=\sqrt{k_n}A_{j,n}(t_n)=0,
 \qquad
 \sqrt{k}\,|A_{j,n}(n/k)|=0\leq B_n(k)
 \quad\text{for every }k\in\mathcal G_n. \tag{9}
\]
The first equality in (9) verifies \(\mathrm{ass:bounded\text{-}bias\text{-}coordinates}\); the inequality verifies \(\mathrm{ass:declared\text{-}bias\text{-}envelope}\), using the declared nonnegativity of \(B_n\).

Equations (2)--(9), together with the structural construction (3), verify every law and sampling property listed in \(\mathrm{def:collision\text{-}model\text{-}class}\).  Hence \((\theta_n^{\mathrm{Par}})_{n\geq1}\in\Theta_{\mathrm{CTC}}(B)\), with \(A_{1,n}=A_{2,n}=0\), for every intermediate sequence \(k_n\) and every declared nonnegative envelope \(B_n\), as claimed.
\end{proof}
\par\smallskip\noindent \textit{Justification.} The transition curve is load-bearing only if its structural witness is inside the tuple class. \textit{Gap.} The anchor uses Pareto innovations in simulations but does not formulate local triangular-array structural membership. \textit{Consumer.} Membership makes the curve an auditable calibration submodel for the marked experiment.

\begin{proposition}[prop:scale-geometry]\label{prop:scale-geometry}
Exactly \(h_n=\kappa_n\lambda_n\). If \(\kappa_n\to0\), every finite local limit lies in \(\mathcal L_0\), so finite \(\lambda\) requires \(h=0\). If \(\kappa_n\to\kappa\in(0,\infty)\), every finite local limit lies in \(\mathcal L_\kappa\), so \(h=\kappa\lambda\). If \(\kappa_n\to\infty\), the finite-dimensional local face is \(\mathcal L_\infty\), so finite \(h\) requires \(\lambda=0\); any nonzero finite limit of \(\lambda_n\) instead satisfies \( |h_n|\to\infty\) and is a separated boundary. For \(k_n=n^\nu\), \(0<\nu<1\), this last regime contains the nonempty band \(k_n^{-1/2}\ll|\Delta_n|\ll1/\log t_n\).
\end{proposition}
\begin{proof}
By \(\mathrm{def{:}collision\text{-}coordinates}\),
\[
 h_n=\sqrt{k_n}\Delta_n
 =\frac{\sqrt{k_n}}{\log t_n}\,\Delta_n\log t_n
 =\kappa_n\lambda_n. \tag{1}
\]
Assumption \(\mathrm{ass{:}intermediate\text{-}sequence}\) implies \(t_n=n/k_n\to\infty\), so \(\log t_n>0\) eventually and every division in (1) is valid.

If \(\kappa_n\to0\) and \(\lambda_n\to\lambda\in\mathbb R\), (1) gives \(h_n\to0\); hence every finite limit has the form \((\lambda,0,b_1,b_2)\in\mathcal L_0\).  If \(\kappa_n\to\kappa\in(0,\infty)\), (1) gives \(h=\kappa\lambda\), exactly the constraint defining \(\mathcal L_\kappa\).  If \(\kappa_n\to\infty\) and \(h_n\to h\in\mathbb R\), then \(\lambda_n=h_n/\kappa_n\to0\), which gives \(\mathcal L_\infty\).  Conversely, if \(\lambda_n\to\lambda\neq0\) in this last regime, then (1) gives \(|h_n|=\kappa_n|\lambda_n|\to\infty\).  These are precisely the three faces in \(\mathrm{def{:}feasible\text{-}local\text{-}spaces}\).

For \(k_n=n^\nu\), \(0<\nu<1\), one has
\[
 t_n=n^{1-\nu},\qquad
 \kappa_n=\frac{n^{\nu/2}}{(1-\nu)\log n}\longrightarrow\infty.
\]
The band is nonempty: for example,
\[
 d_n=\frac{n^{-\nu/4}}{\sqrt{\log n}}
\]
satisfies
\[
 \frac{d_n}{k_n^{-1/2}}
 =\frac{n^{\nu/4}}{\sqrt{\log n}}\longrightarrow\infty,
 \qquad
 d_n\log t_n
 =(1-\nu)n^{-\nu/4}\sqrt{\log n}\longrightarrow0.
\]
Thus \(k_n^{-1/2}\ll d_n\ll1/\log t_n\), proving the last assertion.
\end{proof}
\par\smallskip\noindent \textit{Justification.} The identity \(h_n=\kappa_n\lambda_n\) rules out treating logarithmic mixture drift and root-threshold ordering drift as unconstrained coordinates. \textit{Gap.} The cited marginal root-\(k_n\) theory and fixed-order CTC endpoints do not state this three-face compatibility geometry. \textit{Consumer.} Any repaired local experiment must be indexed on the correct feasible face.

\begin{openendedquestion}[oeq:honest-adaptive-direction-set]\label{oeq:honest-adaptive-direction-set}
Conditionally on the declared envelope \(B_n\), use the finite-grid envelope-conditional handle to construct \((C_n^*)\in\mathcal H_\eta(\Theta_{\mathrm{CTC}}(B))\). For every \(\kappa\in[0,\infty]\) and named compact \(K\subset\mathcal L_\kappa\), give an oracle ambiguity bound over \(K_n(K,\kappa;B)\) in terms of the reported finite grid \(\mathcal G_n\), \(\eta\), \(B_n(k)\), \(s_{\mathrm{CTC},n}(k)\), and \(s_{\mathrm{idx},n}(k)\); require \(\sup_{\theta_n\in K_n(K,\kappa;B)}P_{\theta,n}^{\otimes n}\{C_n^*=\mathcal A_{\mathrm{dir}}\}\to0\) whenever the best grid point has both separations diverging beyond its simultaneous critical value and declared bias envelope. On the reported finite grid \(\mathcal C\), also construct \(C_n^{*,(c)}\in\mathcal H_\eta(\Theta_{\mathrm{CTC}}(B^{(c)}))\) for \(B_n^{(c)}=cB_n\), with \(c\leq c'\) implying \(C_n^{*,(c)}\subseteq C_n^{*,(c')}\). Establish validity at \(\Delta_n=0\) and equivariance under exchanging variable labels.
\end{openendedquestion}
The frozen theorem \(\mathrm{thm{:}constant\text{-}honest\text{-}direction\text{-}set}\) proves that \(C_n^*=\mathcal A_{\mathrm{dir}}\) and \(C_n^{*,(c)}=\mathcal A_{\mathrm{dir}}\) are \(\mathcal T_n\)-measurable, have coverage exactly one at equality and under either ordering, are label-equivariant, and are nested in \(c\). For every eventually nonempty \(K_n(K,\kappa;B)\), their worst-case ambiguity probability equals one.
\par\smallskip\noindent \textit{Justification.} This open question asks whether the constant honest fallback can be replaced by a data-informative envelope-conditional three-valued set on the reported finite threshold grid. \textit{Gap.} The current core has no joint-process calibration, no bias bound for \(\widehat D_n\) derived from the innovation envelope, and no simultaneous coverage proof after threshold selection. The constant full-alphabet set is the strongest proved partial result. \textit{Consumer.} A solution would provide calibrated singleton reports when separation is adequate while retaining an explicit ambiguous report near collision.

\begin{openendedquestion}[oeq:matching-ambiguity-frontier]\label{oeq:matching-ambiguity-frontier}
For each \(\kappa\in[0,\infty]\) and named compact \(K\subset\mathcal L_\kappa\), use the opposite-report tail-rank contiguity handle to determine an explicit \(v_\eta(K)\) such that \[\liminf_{n\to\infty}\inf_{(C_n)\in\mathcal H_\eta(\Theta_{\mathrm{CTC}}(B))}\sup_{\theta_n\in K_n(K,\kappa;B)}P_{\theta,n}^{\otimes n}\{C_n=\mathcal A_{\mathrm{dir}}\}\geq v_\eta(K),\] and prove for the same \(K_n(K,\kappa;B)\), coverage class, loss, tail-rank measurability restriction, and limit order that \[\limsup_{n\to\infty}\sup_{\theta_n\in K_n(K,\kappa;B)}P_{\theta,n}^{\otimes n}\{C_n^*=\mathcal A_{\mathrm{dir}}\}\leq v_\eta(K).\] The converse is expressly conditional on \(\mathcal T_n\) and is not a claim about unrestricted full-distribution causal discovery.
\end{openendedquestion}
Let
\[
U_n(K)=\sup_{\theta_n\in K_n(K,\kappa;B)}
P_{\theta,n}^{\otimes n}\{C_n^*=\mathcal A_{\mathrm{dir}}\}.
\]
Whenever \(K_n(K,\kappa;B)\) is nonempty, \(\mathrm{thm{:}constant\text{-}honest\text{-}direction\text{-}set}\) gives, for every \(\theta_n\) in that set,
\[
P_{\theta,n}^{\otimes n}\{C_n^*=\mathcal A_{\mathrm{dir}}\}=1.
\]
Taking the supremum therefore gives \(U_n(K)=1\), and eventual nonemptiness gives
\[
\limsup_{n\to\infty}U_n(K)=1.
\]
This is the certified ambiguity-one fallback, not a matching upper bound at a proved collision frontier.

When \(K\) contains \((0,0,0,0)\), \(\mathrm{thm{:}four\text{-}coordinate\text{-}ambiguity\text{-}obstruction}\) supplies exactly two experiment certificates: the four displayed local coordinates omit a root-\(k_n\) common-index coordinate, and the summed marginal information omits the nonzero covariance generated by order-\(k_n\) overlap of the two exceedance channels.  These certificates invalidate the proposed four-coordinate summed-marginal route, but they do not construct opposite-orientation arrays with bounded tail-rank likelihood divergence.  Accordingly, no positive lower bound for the minimax ambiguity probability and no value \(v_\eta(K)\) matching the certified upper fallback have been proved.  No stronger assertion is made for compact sets \(K\) that do not contain the origin.
\par\smallskip\noindent \textit{Justification.} This remains the question of whether an augmented tail-rank experiment supports matching upper and lower ambiguity bounds; no numerical frontier is currently claimed. \textit{Gap.} The four-coordinate summed-marginal handle is invalid, and the core contains neither the required paired likelihood and overlap intensity nor a bounded-divergence opposite-orientation construction with verified class membership and separation. \textit{Consumer.} A future solution would distinguish ambiguity forced by tail-rank information from ambiguity caused only by a particular threshold procedure.

\begin{theorem}[thm:joint-collision-experiment-obstruction]\label{thm:joint-collision-experiment-obstruction}
Under the frozen collision model, the requested finite-coordinate uniform local experiment does not exist as stated.  For every regime limit \(\kappa\) and every compact \(K\subset\mathcal L_\kappa\) containing \((0,0,0,0)\), the forward zero-bias exact-Pareto subclass contains two arrays with identical \((\lambda_n,h_n,b_{1,n},b_{2,n})=(0,0,0,0)\) whose restrictions to \(\mathcal T_n\) are asymptotically perfectly distinguishable.  In addition, even at a fixed common index, the two exceedance point measures have order-\(k_n\) overlap, so a joint information form cannot in general be obtained by summing the two marginal intensity informations \(I_{rs}=\sum_{a=1}^2\int\dot\ell_{a,r}\dot\ell_{a,s}\,dM_{\vartheta_0}^{(a)}\).  The same-coordinate Hill certificate alone rules out a uniform local experiment indexed only by \((\lambda,h,b_1,b_2)\), while the positive overlap covariance separately rules out the displayed summed-marginal information form even after a common-index coordinate is added.  A valid formulation must at least add a root-\(k_n\) common-tail-index coordinate and a joint overlap intensity, or an equivalent paired marked point process, before bias drifts and a uniform likelihood remainder can be defined.
\end{theorem}
\begin{proof}
Fix a regime limit \(\kappa\in[0,\infty]\) and a compact set \(K\subset\mathcal L_\kappa\) containing \((0,0,0,0)\).  We construct the first obstruction inside the forward, zero-bias, exact-Pareto subclass named by \(\mathrm{def{:}pareto\text{-}witness}\).  Set \(\beta_n=\beta_0\), \(c_{1,n}=c_0\), \(c_{2,n}=1\), and \(A_{1,n}=A_{2,n}=0\).  In array \(s=0\), set
\[
 a_{0,n}=\alpha_{1,n}=\alpha_{2,n}=\alpha_0,
\]
whereas in array \(s=1\), set
\[
 a_{1,n}=\alpha_{1,n}=\alpha_{2,n}=\alpha_0+k_n^{-1/4}.
\]
Because \(k_n\to\infty\) by \(\mathrm{ass{:}intermediate\text{-}sequence}\), both common indices converge to \(\alpha_0\).  In both arrays \(\Delta_n=0\), so \(\mathrm{def{:}collision\text{-}coordinates}\) gives
\[
 (\lambda_n,h_n,b_{1,n},b_{2,n})=(0,0,0,0). \tag{1}
\]
Thus (1) lies in \(K\).  The exact-Pareto quantile is \(U_{j,n}(t)=(c_{j,n}t)^{1/a_{s,n}}\), whence the exact branch of \(\mathrm{ass{:}uniform\text{-}second\text{-}order\text{-}variation}\) holds with \(A_{j,n}=0\), and \(\mathrm{ass{:}uniform\text{-}tail\text{-}normalization}\) holds with equality.  The constants obey \(c_{1,n}/c_{2,n}=c_0\), the innovations are independent and continuous, \(\beta_n\to\beta_0>0\), and the declared nonnegative envelope contains the zero auxiliaries.  Consequently every condition in \(\mathrm{def{:}collision\text{-}model\text{-}class}\) is satisfied.  In the regime for which \(\kappa_n\to\kappa\), both arrays therefore supply members of \(K_n(K,\kappa;B)\).

The marginal upper-order-statistic arrays are measurable with respect to \(\mathcal T_n\) by \(\mathrm{def{:}tail\text{-}rank\text{-}experiment}\).  From the first marginal define
\[
 H_n=\frac1{k_n}\sum_{j=1}^{k_n}
 \log\!\left(\frac{X_{(n-j+1),1,n}}{X_{(n-k_n),1,n}}\right).
\]
Under common Pareto index \(a>0\), the centered log observations are independent exponential variables with rate \(a\).  If \(S_{j,n}\) is the spacing between the \(j\)-th and \((j+1)\)-st largest log observations, the joint exponential density and a triangular change of variables show that \(S_{1,n},\ldots,S_{k_n,n}\) are independent and \(S_{j,n}\) has exponential rate \(ja\).  Expanding the preceding sum gives
\[
 H_n=\frac1{k_n}\sum_{j=1}^{k_n}jS_{j,n}
 \ \stackrel{d}{=}\ \frac1{k_n}\sum_{j=1}^{k_n}E_{j,n},
 \qquad E_{1,n},\ldots,E_{k_n,n}\overset{\mathrm{iid}}{\sim}\operatorname{Exp}(a).
\]
Direct integration of the exponential density yields
\[
 E_aH_n=\frac1a,\qquad
 \operatorname{Var}_a(H_n)=\frac1{k_na^2}. \tag{2}
\]
For the two constructed arrays, put \(\mu_{s,n}=1/a_{s,n}\).  Then
\[
 d_n:=|\mu_{1,n}-\mu_{0,n}|
 =\frac{k_n^{-1/4}}{\alpha_0(\alpha_0+k_n^{-1/4})}
 \asymp k_n^{-1/4}. \tag{3}
\]
Use the midpoint test that selects array \(1\) when \(H_n<(\mu_{0,n}+\mu_{1,n})/2\).  For either array, the error event implies \(|H_n-\mu_{s,n}|\geq d_n/2\).  Since
\[
 E_s[(H_n-\mu_{s,n})^2]
 \geq \frac{d_n^2}{4}
 P_s\{|H_n-\mu_{s,n}|\geq d_n/2\},
\]
equations (2)--(3) imply that the sum of the two error probabilities is \(O(k_n^{-1/2})\to0\).  Hence the total-variation distance between the two restrictions to \(\mathcal T_n\) is at least one minus this error sum and therefore tends to one.  Nevertheless, their four coordinates are identical by (1).  Moreover, the omitted common-index coordinate
\[
 \sqrt{k_n}\left\{\frac{\alpha_{1,n}+\alpha_{2,n}}2-\alpha_0\right\}
\]
is zero in the first array and equals \(k_n^{1/4}\to\infty\) in the second, while (2) identifies the root-\(k_n\) stochastic scale.  Thus no likelihood expansion indexed uniformly only by the four coordinates in (1) can represent these restrictions.

We next prove the overlap obstruction.  Keep array \(0\), and let
\[
 Y_{a,n}=N_n^{(a)}((1,\infty)\times[0,1]),
 \qquad a\in\{1,2\}.
\]
The first marginal is exact Pareto, so
\[
 q_{1,n}=(c_0t_n)^{1/\alpha_0}. \tag{4}
\]
For completeness, write \(A=\beta_0\epsilon_{1,n}\), \(E=\epsilon_{2,n}\), and fix \(\delta\in(0,1/2)\).  Independence and nonnegativity give
\[
 P(A>x)+P(E>x)-P(A>x)P(E>x)
 \leq P(A+E>x),
\]
and the deterministic inclusion
\[
 \{A+E>x\}\subset
 \{A>(1-\delta)x\}\cup\{E>(1-\delta)x\}
 \cup\{A>\delta x,E>\delta x\}
\]
gives the reverse upper bound.  Multiplication by \(x^{\alpha_0}\), passage to the limit, and then \(\delta\downarrow0\) therefore yield
\[
 P(\beta_0\epsilon_{1,n}+\epsilon_{2,n}>x)
 =(\beta_0^{\alpha_0}c_0+1)x^{-\alpha_0}\{1+o(1)\}. \tag{5}
\]
The convolution has a continuous, eventually strictly decreasing tail.  Since its \((1-1/t_n)\)-quantile is \(q_{2,n}\), (5) and \(t_n\to\infty\) imply
\[
 q_{2,n}^{\alpha_0}/t_n\longrightarrow
 \beta_0^{\alpha_0}c_0+1.
\]
Combining this relation with (4) gives
\[
 \frac{q_{2,n}}{q_{1,n}}
 \longrightarrow
 \left(\frac{\beta_0^{\alpha_0}c_0+1}{c_0}\right)^{1/\alpha_0},
 \qquad
 \frac{q_{2,n}}{\beta_0q_{1,n}}
 \longrightarrow
 \left(1+\frac1{\beta_0^{\alpha_0}c_0}\right)^{1/\alpha_0}>1. \tag{6}
\]
For all sufficiently large \(n\), (6) makes
\[
 \{\epsilon_{1,n}>q_{2,n}/\beta_0\}
 \subseteq
 \{X_{1,n}>q_{1,n},\ X_{2,n}>q_{2,n}\}. \tag{7}
\]
Continuity at the marginal quantiles gives \(P(X_{a,n}>q_{a,n})=1/t_n\).  Independence across observations under \(\mathrm{ass{:}iid\text{-}triangular\text{-}sampling}\), followed by (5)--(7), yields
\[
\begin{aligned}
 \frac1{k_n}\operatorname{Cov}(Y_{1,n},Y_{2,n})
 &=\frac n{k_n}\left[
 P(X_{1,n}>q_{1,n},X_{2,n}>q_{2,n})-t_n^{-2}\right]\\
 &\geq t_nP(\epsilon_{1,n}>q_{2,n}/\beta_0)-t_n^{-1}\\
 &\longrightarrow
 \frac{\beta_0^{\alpha_0}c_0}{1+\beta_0^{\alpha_0}c_0}>0.
\end{aligned}\tag{8}
\]
Thus the two compensated exceedance measures have a nonzero limiting cross-covariance already for the constant test functions defining \(Y_{1,n}\) and \(Y_{2,n}\).  The separate normalized marginal intensities \(M_{n,\theta}^{(1)}\) and \(M_{n,\theta}^{(2)}\) do not encode this overlap measure, and the candidate form
\[
 I_{rs}=\sum_{a=1}^2\int
 \dot\ell_{a,r}(z)\dot\ell_{a,s}(z)\,dM_{\vartheta_0}^{(a)}(z)
\]
in \(\mathrm{def{:}joint\text{-}experiment\text{-}handle}\) has no cross term.  Equations (1)--(8) prove the two asserted obstructions.  They also show why a common-index coordinate at root-\(k_n\) scale and a joint overlap intensity, or an equivalent paired marked process, are necessary before bias drifts and a uniform likelihood remainder can be specified.
\end{proof}
\par\smallskip\noindent \textit{Justification.} The two certificates are logically separate: the Hill statistic rules out four-coordinate uniformity, while positive overlap covariance rules out the displayed summed-marginal information formula even after adding a common-index coordinate. \textit{Gap.} Marginal tail-index local asymptotic normality and one-process tail-copula limits do not provide the omitted common-index coordinate or paired overlap intensity. \textit{Consumer.} It determines the honest inference fallback and specifies the augmentation required before nonconstant calibration can be attempted.

\begin{theorem}[thm:constant-honest-direction-set]\label{thm:constant-honest-direction-set}
Let \(C_n^0\) be the constant sample map
\[
 C_n^*(z_1,\ldots,z_n)=C_n^0(z_1,\ldots,z_n)=\mathcal A_{\mathrm{dir}}
 \qquad\text{for every }(z_1,\ldots,z_n)\in(\mathbb R_+^2)^n,
\]
and, for every \(c\in\mathcal C\), let \(C_n^{*,(c)}=\mathcal A_{\mathrm{dir}}\).  For every \(\eta\in(0,1/2)\),
\[
 (C_n^*)\in\mathcal H_\eta(\Theta_{\mathrm{CTC}}(B)),
 \qquad
 (C_n^{*,(c)})\in\mathcal H_\eta(\Theta_{\mathrm{CTC}}(B^{(c)}))
 \quad(c\in\mathcal C).
\]
The coverage probability is exactly one for every model member, including every member with \(\Delta_n=0\).  The construction is equivariant under exchanging the two variable labels, and
\[
 c\leq c'\quad\Longrightarrow\quad
 C_n^{*,(c)}=C_n^{*,(c')}=\mathcal A_{\mathrm{dir}},
\]
so the sensitivity family is nested.  Moreover,
\[
 P_{\theta,n}^{\otimes n}\{C_n^*=\mathcal A_{\mathrm{dir}}\}=1
\]
for every \(\theta_n\); hence its worst-case ambiguity probability over every nonempty \(K_n(K,\kappa;B)\) is exactly one.  This is the strongest direction-set conclusion established from the frozen calibration handle: the declared envelope controls \(\sqrt{k}\lvert A_{j,n}(n/k)\rvert\), not the bias of \(\widehat D_n\), while the proposed four-coordinate summed-marginal experiment omits a common-index coordinate and the order-\(k_n\) overlap of the two exceedance channels.  No vanishing-ambiguity conclusion is asserted.
\end{theorem}
\begin{proof}
Fix \(n\), and write \(\Omega_n=(\mathbb R_+^2)^n\).  The displayed construction makes \(C_n^*\) the constant map from \(\Omega_n\) to \(\mathcal A_{\mathrm{dir}}\).  The inverse image of any subset of its finite range is either \(\varnothing\) or \(\Omega_n\); both belong to \(\mathcal T_n\).  Thus \(C_n^*\) is \(\mathcal T_n\)-measurable.  The same argument applies to every \(C_n^{*,(c)}\), \(c\in\mathcal C\).

Let \(\theta=(\theta_m)_{m\geq1}\in\Theta_{\mathrm{CTC}}(B)\).  By the orientation projection and the definition of the alphabet,
\[
 o(\theta_n)=o_n\in\{\mathsf{forward},\mathsf{reverse}\}
 =\mathcal A_{\mathrm{dir}}=C_n^*.
\]
Consequently
\[
 P_{\theta,n}^{\otimes n}\{o(\theta_n)\in C_n^*\}=1. \tag{1}
\]
If the model class is nonempty, its infimum in (1) is one; if it is empty, the universal coverage requirement is vacuous.  Hence in either case
\[
 \liminf_{n\to\infty}\inf_{\theta\in\Theta_{\mathrm{CTC}}(B)}
 P_{\theta,n}^{\otimes n}\{o(\theta_n)\in C_n^*\}
 \geq1\geq1-\eta,
\]
where \(\eta\in(0,1/2)\).  Together with measurability and the stated range, \(\mathrm{def{:}honest\text{-}direction\text{-}class}\) gives
\[
 (C_n^*)\in\mathcal H_\eta(\Theta_{\mathrm{CTC}}(B)). \tag{2}
\]
Replacing \(B\) by \(B^{(c)}\) does not change the pointwise inclusion of the orientation in \(\mathcal A_{\mathrm{dir}}\).  Repeating (1)--(2) proves
\[
 (C_n^{*,(c)})\in\mathcal H_\eta(\Theta_{\mathrm{CTC}}(B^{(c)}))
 \qquad(c\in\mathcal C).
\]
No step used the sign or magnitude of \(\Delta_n\), so coverage remains one at \(\Delta_n=0\) and under either strict ordering.

Let \(\sigma\) exchange both observed-variable labels and the two orientation labels.  Since \(\sigma(\mathcal A_{\mathrm{dir}})=\mathcal A_{\mathrm{dir}}\),
\[
 C_n^*(\sigma z)=\mathcal A_{\mathrm{dir}}
 =\sigma(\mathcal A_{\mathrm{dir}})=\sigma C_n^*(z),
\]
which proves equivariance.  If \(c\leq c'\), then
\[
 C_n^{*,(c)}=\mathcal A_{\mathrm{dir}}=C_n^{*,(c')},
\]
so nesting holds with equality.  Finally,
\[
 \{C_n^*=\mathcal A_{\mathrm{dir}}\}=\Omega_n,
 \qquad
 P_{\theta,n}^{\otimes n}\{C_n^*=\mathcal A_{\mathrm{dir}}\}=1. \tag{3}
\]
Taking the supremum of (3) over any nonempty \(K_n(K,\kappa;B)\) leaves the value one.

It remains to justify the stated scope.  Assumption \(\mathrm{ass{:}declared\text{-}bias\text{-}envelope}\) controls only
\[
 \sqrt{k}\,|A_{j,n}(n/k)|\leq B_n(k)
\]
and provides no bound for the bias of \(\widehat D_n\).  Moreover, \(\mathrm{thm{:}joint\text{-}collision\text{-}experiment\text{-}obstruction}\) proves that the four collision coordinates omit a root-\(k_n\) common-index coordinate and that the two exceedance channels have an order-\(k_n\) overlap absent from the summed marginal information.  The frozen premises therefore provide neither the contrast-bias control nor the joint critical value named by \(\mathrm{def{:}adaptive\text{-}direction\text{-}handle}\).  Equations (1)--(3) establish the certified exactly honest fallback and its ambiguity value, without asserting vanishing ambiguity or impossibility of every nonconstant procedure.
\end{proof}
\par\smallskip\noindent \textit{Justification.} The constant sample map returning \(\mathcal A_{\mathrm{dir}}\) is measurable and contains every structural orientation pointwise, so its coverage is exactly one for the full class, at equality, under either ordering, and for every reported envelope multiplier. \textit{Gap.} The frozen premises do not provide a paired marked-process critical value or a bound converting the declared second-order envelope into bias control for \(\widehat D_n\), so no nonconstant calibrated direction rule is promoted. \textit{Consumer.} Climate and finance analyses can report the full direction alphabet as an exactly covered fallback set; no calibrated arrow or vanishing-ambiguity claim accompanies it.

\begin{theorem}[thm:class-level-collision-curve]\label{thm:class-level-collision-curve}
For every triangular array \((\theta_n)_{n\geq1}\in\Theta_{\mathrm{CTC}}(B)\) such that \(o_n=\mathsf{forward}\) eventually and
\[
\xi_n=\Delta_n\log U_{2,n}(t_n)\longrightarrow\xi\in\mathbb R,
\]
put
\[
R_\xi=\beta_0^{\alpha_0}c_0\exp(-\xi).
\]
Then
\[
\Gamma_{12,n}(t_n)\longrightarrow1,
\qquad
\Gamma_{21,n}(t_n)\longrightarrow
\frac12+\frac{R_\xi}{2(1+R_\xi)},
\]
and
\[
D_n(t_n)\longrightarrow\frac{1}{2(1+R_\xi)}.
\]
These are sequential conclusions for each qualifying array; no additional uniformity over arrays or class-uniform remainder is asserted. In particular, every eventually-forward array in \(\Theta_{\mathrm{CTC}}(B)\) satisfying \(\Delta_n\log c_{2,n}\to0\) and \(\lambda_n=\Delta_n\log t_n\to\lambda\in\mathbb R\) obeys
\[
\xi_n\longrightarrow\xi=\frac{\lambda}{\alpha_0},
\qquad
R_\xi=R(\lambda)=\beta_0^{\alpha_0}c_0\exp(-\lambda/\alpha_0),
\]
so the preceding limits hold with \(R_\xi=R(\lambda)\), and
\[
D_n(t_n)\longrightarrow g_{\alpha_0,\beta_0,c_0}(\lambda)
=\frac{1}{2\{1+R(\lambda)\}}.
\]
\end{theorem}
\begin{proof}
Fix one triangular array in \(\Theta_{\mathrm{CTC}}(B)\) satisfying the intrinsic-coordinate hypotheses. All limits below are along this fixed array. By eventual forward orientation and \(\mathrm{ass{:}linear\text{-}scm}\),
\[
X_{1,n}=\epsilon_{1,n},
\qquad
X_{2,n}=\beta_n\epsilon_{1,n}+\epsilon_{2,n}.
\tag{1}
\]
For brevity, write
\[
U_{j,n}=U_{j,n}(t_n),
\qquad
a_n=\beta_nU_{1,n},
\qquad
b_n=U_{2,n},
\qquad
M_n=\max(a_n,b_n).
\tag{2}
\]
The definitions of \(c_{1,n}\) and \(c_{2,n}\) give the exact identities
\[
U_{j,n}^{\alpha_{j,n}}=c_{j,n}t_n,
\qquad j\in\{1,2\}.
\tag{3}
\]
Eliminating \(t_n\) from the two identities in (3) yields
\[
\log\frac{U_{1,n}}{U_{2,n}}
=\frac{\log(c_{1,n}/c_{2,n})}{\alpha_{1,n}}
-\frac{\Delta_n}{\alpha_{1,n}}\log U_{2,n}.
\tag{4}
\]
By \(\mathrm{ass{:}tail\text{-}constant\text{-}balance}\), \(\log(c_{1,n}/c_{2,n})\to\log c_0\). By \(\mathrm{ass{:}local\text{-}tail\text{-}collision}\), \(\alpha_{1,n}\to\alpha_0>0\) and \(\Delta_n\to0\). The hypothesis \(\xi_n=\Delta_n\log U_{2,n}\to\xi\) therefore turns (4) into
\[
\log\frac{U_{1,n}}{U_{2,n}}
\longrightarrow \frac{\log c_0-\xi}{\alpha_0}.
\tag{5}
\]
Together with \(\beta_n\to\beta_0>0\) from \(\mathrm{ass{:}coefficient\text{-}limit}\), (5) supplies a constant \(K<\infty\) such that, eventually,
\[
K^{-1}\leq\frac{a_n}{b_n}\leq K.
\tag{6}
\]

Let \(q_{2,n}=F_{2,n}^{\leftarrow}(1-1/t_n)\). Assumption \(\mathrm{ass{:}continuous\text{-}margins}\) implies
\[
P_{\theta,n}(X_{2,n}>q_{2,n})=t_n^{-1},
\tag{7}
\]
and the conditioning events \(\{F_{2,n}(X_{2,n})>1-1/t_n\}\) and \(\{X_{2,n}>q_{2,n}\}\) agree up to a null set. By \(\mathrm{ass{:}local\text{-}tail\text{-}collision}\), there are constants \(0<\underline\alpha<\overline\alpha<\infty\) such that, eventually,
\[
\underline\alpha\leq\alpha_{j,n}\leq\overline\alpha,
\qquad j\in\{1,2\}.
\tag{8}
\]
Fix \(d\in(0,1)\). If \(M_n=a_n\), then \(\mathrm{ass{:}uniform\text{-}tail\text{-}normalization}\), applied at multiplier \(d\), gives
\[
t_nP_{\theta,n}(\beta_n\epsilon_{1,n}>dM_n)
=d^{-\alpha_{1,n}}\{1+o(1)\}>1
\tag{9}
\]
eventually. If \(M_n=b_n\), the same conclusion follows with \(\epsilon_{2,n}\) and \(\alpha_{2,n}\). Positivity in (1) and (7) then force \(q_{2,n}\geq dM_n\).

For the converse comparison, choose a fixed \(C>2\) such that \(2(C/2)^{-\underline\alpha}<1/2\). By (6), the multipliers \(CM_n/(2a_n)\) and \(CM_n/(2b_n)\) range in a fixed compact subset of \((0,\infty)\), and each is at least \(C/2\). The union bound, (8), and \(\mathrm{ass{:}uniform\text{-}tail\text{-}normalization}\) give
\[
\begin{aligned}
P_{\theta,n}(X_{2,n}>CM_n)
&\leq P_{\theta,n}(\beta_n\epsilon_{1,n}>CM_n/2)
+P_{\theta,n}(\epsilon_{2,n}>CM_n/2),\\
t_nP_{\theta,n}(X_{2,n}>CM_n)
&\leq 2(C/2)^{-\underline\alpha}\{1+o(1)\}<1
\end{aligned}
\tag{10}
\]
eventually. Equation (7) therefore forces \(q_{2,n}\leq CM_n\). We have proved the compact quantile sandwich
\[
dM_n\leq q_{2,n}\leq CM_n.
\tag{11}
\]
In particular, \(q_{2,n}/a_n\) and \(q_{2,n}/b_n\) are bounded above and bounded away from zero.

Define
\[
p_{1,n}=P_{\theta,n}(\beta_n\epsilon_{1,n}>q_{2,n}),
\qquad
p_{2,n}=P_{\theta,n}(\epsilon_{2,n}>q_{2,n}).
\tag{12}
\]
Using (3), (11), and \(\mathrm{ass{:}uniform\text{-}tail\text{-}normalization}\) at the two compact multiplier sequences gives
\[
\begin{aligned}
p_{1,n}
&=t_n^{-1}\left(\frac{q_{2,n}}{a_n}\right)^{-\alpha_{1,n}}\{1+o(1)\},\\
p_{2,n}
&=t_n^{-1}\left(\frac{q_{2,n}}{b_n}\right)^{-\alpha_{2,n}}\{1+o(1)\}.
\end{aligned}
\tag{13}
\]
Thus \(p_{1,n}\asymp t_n^{-1}\) and \(p_{2,n}\asymp t_n^{-1}\).

Write \(Y_n=\beta_n\epsilon_{1,n}\) and \(W_n=\epsilon_{2,n}\). For any fixed \(\delta\in(0,1/2)\), positivity gives the event bounds
\[
\begin{aligned}
p_{1,n}+p_{2,n}-p_{1,n}p_{2,n}
&\leq P_{\theta,n}(Y_n+W_n>q_{2,n})\\
&\leq P_{\theta,n}(Y_n>(1-\delta)q_{2,n})
+P_{\theta,n}(W_n>(1-\delta)q_{2,n})\\
&\quad+P_{\theta,n}(Y_n>\delta q_{2,n})
P_{\theta,n}(W_n>\delta q_{2,n}).
\end{aligned}
\tag{14}
\]
Here both product identities use \(\mathrm{ass{:}innovation\text{-}independence}\). By (11), all thresholds in (14), divided by the corresponding \(U_{j,n}\), remain in fixed compact subsets of \((0,\infty)\). Hence \(\mathrm{ass{:}uniform\text{-}tail\text{-}normalization}\) and (13) imply
\[
\frac{P_{\theta,n}(Y_n>(1-\delta)q_{2,n})}{p_{1,n}}
=(1-\delta)^{-\alpha_{1,n}}+o(1),
\tag{15}
\]
with the analogous identity for \(W_n\), while the product in the last line of (14) is \(O(t_n^{-2})\). Assumption \(\mathrm{ass{:}intermediate\text{-}sequence}\) gives \(t_n=n/k_n\to\infty\), so that product and \(p_{1,n}p_{2,n}\) are \(o(p_{1,n}+p_{2,n})\). Dividing (14) by \(p_{1,n}+p_{2,n}\), taking lower and upper limits, and using \(\alpha_{j,n}\to\alpha_0\), we obtain, for every fixed \(\delta\),
\[
1\leq\liminf_{n\to\infty}
\frac{P_{\theta,n}(Y_n+W_n>q_{2,n})}{p_{1,n}+p_{2,n}}
\leq\limsup_{n\to\infty}
\frac{P_{\theta,n}(Y_n+W_n>q_{2,n})}{p_{1,n}+p_{2,n}}
\leq(1-\delta)^{-\alpha_0}.
\]
Letting \(\delta\downarrow0\) proves the ratio-free one-large-jump expansion
\[
P_{\theta,n}(X_{2,n}>q_{2,n})
=p_{1,n}+p_{2,n}+o(p_{1,n}+p_{2,n}).
\tag{16}
\]

Put \(H_n=F_{1,n}(X_{1,n})=F_{1,n}(\epsilon_{1,n})\). For \(v\in(0,1)\), let \(x_v=F_{1,n}^{\leftarrow}(v)\). Continuity gives \(F_{1,n}(x_v)=v\), and \(\{F_{1,n}(X_{1,n})\leq v\}\) agrees almost surely with \(\{X_{1,n}\leq x_v\}\). Hence \(P_{\theta,n}(H_n\leq v)=v\), so \(H_n\) is uniform on \((0,1)\). Therefore
\[
E_{P_{\theta,n}}(H_n)=\frac12.
\tag{17}
\]
On \(\{Y_n>q_{2,n}\}\), monotonicity of \(F_{1,n}\) gives
\[
0\leq1-H_n
\leq P_{\theta,n}(\epsilon_{1,n}>q_{2,n}/\beta_n)=p_{1,n}.
\tag{18}
\]
Consequently,
\[
E_{P_{\theta,n}}[H_n\mathbf 1\{Y_n>q_{2,n}\}]
=p_{1,n}+O(p_{1,n}^2),
\qquad
E_{P_{\theta,n}}[H_n\mathbf 1\{W_n>q_{2,n}\}]
=\frac12p_{2,n},
\tag{19}
\]
where the second equality uses \(\mathrm{ass{:}innovation\text{-}independence}\) and (17). The intersection \(\{Y_n>q_{2,n},W_n>q_{2,n}\}\) has probability \(p_{1,n}p_{2,n}=o(p_{1,n}+p_{2,n})\). By (16), the part of \(\{Y_n+W_n>q_{2,n}\}\) outside the union of those two events also has probability \(o(p_{1,n}+p_{2,n})\). Since \(0\leq H_n\leq1\), equations (16)--(19) give
\[
\Gamma_{21,n}(t_n)
=\frac{p_{1,n}+p_{2,n}/2}{p_{1,n}+p_{2,n}}+o(1)
=\frac12+\frac12\frac{p_{1,n}/p_{2,n}}{1+p_{1,n}/p_{2,n}}+o(1).
\tag{20}
\]

Equation (13) also gives the asymptotic ratio formula
\[
\frac{p_{1,n}}{p_{2,n}}
=\beta_n^{\alpha_{1,n}}\frac{c_{1,n}}{c_{2,n}}
q_{2,n}^{-\Delta_n}\{1+o(1)\}.
\tag{21}
\]
Because (5) makes \(\log(U_{1,n}/U_{2,n})\) bounded and \(\beta_n\to\beta_0>0\), while \(\Delta_n\to0\),
\[
\begin{aligned}
\Delta_n\log b_n&=\xi_n\longrightarrow\xi,\\
\Delta_n\log a_n
&=\xi_n+\Delta_n\log\frac{U_{1,n}}{U_{2,n}}
+\Delta_n\log\beta_n\longrightarrow\xi.
\end{aligned}
\tag{22}
\]
Thus \(\Delta_n\log M_n\to\xi\), regardless of which term realizes the maximum. Equation (11) gives \(d\leq q_{2,n}/M_n\leq C\); multiplying the bounded logarithm of this ratio by \(\Delta_n\to0\) yields
\[
\Delta_n\log q_{2,n}\longrightarrow\xi.
\tag{23}
\]
Substituting (23), \(\beta_n\to\beta_0\), \(\alpha_{1,n}\to\alpha_0\), and \(c_{1,n}/c_{2,n}\to c_0\) into (21) proves
\[
\frac{p_{1,n}}{p_{2,n}}
\longrightarrow
\beta_0^{\alpha_0}c_0e^{-\xi}=R_\xi.
\tag{24}
\]
Equations (20) and (24) establish the asserted limit for \(\Gamma_{21,n}(t_n)\).

It remains to evaluate the other rank coefficient. Under (1), \(F_{1,n}\) is the continuous distribution function of \(\epsilon_{1,n}\), so its \((1-1/t_n)\)-quantile is \(U_{1,n}\). On \(\{X_{1,n}>U_{1,n}\}\), positivity and (1) imply \(X_{2,n}>\beta_nU_{1,n}\). Monotonicity of \(F_{2,n}\) therefore gives
\[
\begin{aligned}
0\leq1-\Gamma_{12,n}(t_n)
&=E_{P_{\theta,n}}[1-F_{2,n}(X_{2,n})\mid X_{1,n}>U_{1,n}]\\
&\leq P_{\theta,n}(X_{2,n}>\beta_nU_{1,n}).
\end{aligned}
\tag{25}
\]
The two multipliers in the following union bound are \(1/2\) relative to \(U_{1,n}\) and \(a_n/(2b_n)\) relative to \(U_{2,n}\); by (6), both remain in a fixed compact subset of \((0,\infty)\). Thus \(\mathrm{ass{:}uniform\text{-}tail\text{-}normalization}\) gives
\[
\begin{aligned}
P_{\theta,n}(X_{2,n}>\beta_nU_{1,n})
&\leq P_{\theta,n}(\epsilon_{1,n}>U_{1,n}/2)
+P_{\theta,n}(\epsilon_{2,n}>\beta_nU_{1,n}/2)\\
&=O(t_n^{-1})=o(1),
\end{aligned}
\tag{26}
\]
where the last equality uses \(t_n\to\infty\). Equations (25)--(26) prove \(\Gamma_{12,n}(t_n)\to1\). By \(\mathrm{def{:}finite\text{-}threshold\text{-}ctc}\), (20), and (24),
\[
D_n(t_n)
\longrightarrow
1-\left\{\frac12+\frac{R_\xi}{2(1+R_\xi)}\right\}
=\frac{1}{2(1+R_\xi)}.
\tag{27}
\]

Finally, consider an eventually-forward array satisfying \(\Delta_n\log c_{2,n}\to0\) and \(\lambda_n=\Delta_n\log t_n\to\lambda\in\mathbb R\). The \(j=2\) identity in (3) implies
\[
\alpha_{2,n}\xi_n
=\Delta_n\log c_{2,n}+\Delta_n\log t_n
=\Delta_n\log c_{2,n}+\lambda_n
\longrightarrow\lambda.
\tag{28}
\]
Since \(\alpha_{2,n}\to\alpha_0>0\), division in (28) is valid eventually and gives \(\xi_n\to\lambda/\alpha_0\). Consequently,
\[
R_\xi
=\beta_0^{\alpha_0}c_0e^{-\lambda/\alpha_0}
=R(\lambda),
\tag{29}
\]
and (20), (24), and (27) become
\[
\Gamma_{12,n}(t_n)\to1,
\qquad
\Gamma_{21,n}(t_n)\to\frac12+\frac{R(\lambda)}{2\{1+R(\lambda)\}},
\qquad
D_n(t_n)\to g_{\alpha_0,\beta_0,c_0}(\lambda).
\]
This proves the intrinsic-coordinate theorem and its declared \(\lambda_n\)-coordinate corollary, with sequential-per-array scope throughout.
\end{proof}
\par\smallskip\noindent \textit{Justification.} The intrinsic threshold coordinate yields a strictly broader class-level population collision curve, while the declared lambda curve follows as an exact corollary. \textit{Gap.} Published fixed-index regimes do not supply this sequential intrinsic-coordinate map or its reduction to the local lambda scale, and no published inference result is transferred. \textit{Consumer.} It gives the population geometry that the joint local experiment, reverse-orientation reduction, and any honest direction procedure must target.

\begin{theorem}[thm:reverse-collision-curve]\label{thm:reverse-collision-curve}
For every triangular array \((\theta_n)_{n\geq1}\in\Theta_{\mathrm{CTC}}(B)\) such that \(o_n=\mathsf{reverse}\) eventually and
\[
\xi_{\mathrm{rev},n}=-\Delta_n\log U_{1,n}(t_n)
\longrightarrow\xi_{\mathrm{rev}}\in\mathbb R,
\]
put
\[
R_{\mathrm{rev},\xi}
=\beta_0^{\alpha_0}c_0^{-1}\exp(-\xi_{\mathrm{rev}}).
\]
Then
\[
\Gamma_{21,n}(t_n)\longrightarrow1,
\qquad
\Gamma_{12,n}(t_n)\longrightarrow
\frac12+\frac{R_{\mathrm{rev},\xi}}{2(1+R_{\mathrm{rev},\xi})},
\]
and
\[
D_n(t_n)\longrightarrow-\frac{1}{2(1+R_{\mathrm{rev},\xi})}.
\]
These are sequential conclusions for each qualifying reverse array; no additional uniformity over arrays or class-uniform remainder is asserted. In particular, every eventually-reverse array in \(\Theta_{\mathrm{CTC}}(B)\) satisfying \(\Delta_n\log c_{2,n}\to0\) and \(\lambda_n=\Delta_n\log t_n\to\lambda\in\mathbb R\) obeys
\[
\xi_{\mathrm{rev},n}\longrightarrow\xi_{\mathrm{rev}}=-\frac{\lambda}{\alpha_0}.
\]
Writing
\[
R_{\mathrm{rev}}(\lambda)
=\beta_0^{\alpha_0}c_0^{-1}\exp(\lambda/\alpha_0),
\]
we have \(R_{\mathrm{rev},\xi}=R_{\mathrm{rev}}(\lambda)\), the preceding limits hold with this ratio, and
\[
D_n(t_n)\longrightarrow
-g_{\alpha_0,\beta_0,1/c_0}(-\lambda).
\]
On the reverse scaled exact-Pareto subclass with independent innovations satisfying \(Q_{j,n}\{\epsilon>x\}=c_{j,n}x^{-\alpha_{j,n}}\) above their support endpoints and structural equations \(X_{2,n}=\epsilon_{2,n}\), \(X_{1,n}=\beta_nX_{2,n}+\epsilon_{1,n}\), the condition \(\Delta_n\log c_{2,n}\to0\) and \(\lambda_n\to\lambda\in[-\infty,+\infty]\) give the extended endpoints
\[
(\Gamma_{21,n}(t_n),\Gamma_{12,n}(t_n),D_n(t_n))\longrightarrow
\begin{cases}
(1,\frac12,-\frac12), & \lambda=-\infty,\\
\left(1,\frac12+\frac{R_{\mathrm{rev}}(\lambda)}{2\{1+R_{\mathrm{rev}}(\lambda)\}},-g_{\alpha_0,\beta_0,1/c_0}(-\lambda)\right), & \lambda\in\mathbb R,\\
(1,1,0), & \lambda=+\infty.
\end{cases}
\]
\end{theorem}
\begin{proof}
Fix one triangular array satisfying the intrinsic-coordinate hypotheses. All limits until the exact-Pareto paragraph are along this fixed array. By eventual reverse orientation and \(\mathrm{ass{:}linear\text{-}scm}\),
\[
X_{2,n}=\epsilon_{2,n},
\qquad
X_{1,n}=\beta_n\epsilon_{2,n}+\epsilon_{1,n}.
\tag{1}
\]
For brevity, put
\[
u_{j,n}=U_{j,n}(t_n),
\qquad
a_n=u_{1,n},
\qquad
b_n=\beta_n u_{2,n},
\qquad
M_n=\max(a_n,b_n).
\tag{2}
\]
The definition of the threshold-normalized constants gives the exact identities
\[
u_{j,n}^{\alpha_{j,n}}=c_{j,n}t_n,
\qquad j\in\{1,2\}.
\tag{3}
\]
Subtracting the logarithms of the two identities in (3) and using \(\Delta_n=\alpha_{1,n}-\alpha_{2,n}\) yields
\[
\log\frac{u_{2,n}}{u_{1,n}}
=\frac{-\log(c_{1,n}/c_{2,n})+\Delta_n\log u_{1,n}}
{\alpha_{2,n}}.
\tag{4}
\]
Assumption \(\mathrm{ass{:}tail\text{-}constant\text{-}balance}\) implies
\(\log(c_{1,n}/c_{2,n})\to\log c_0\). Assumption
\(\mathrm{ass{:}local\text{-}tail\text{-}collision}\) implies
\(\alpha_{2,n}\to\alpha_0>0\) and \(\Delta_n\to0\). The intrinsic-coordinate hypothesis therefore turns (4) into
\[
\log\frac{u_{2,n}}{u_{1,n}}
\longrightarrow
\frac{-\log c_0-\xi_{\mathrm{rev}}}{\alpha_0}.
\tag{5}
\]
Together with \(\beta_n\to\beta_0>0\) from
\(\mathrm{ass{:}coefficient\text{-}limit}\), equation (5) supplies a finite
constant \(K\geq1\) such that, eventually,
\[
K^{-1}\leq\frac{b_n}{a_n}\leq K.
\tag{6}
\]
All logarithms above are well defined because the quantiles, tail constants, and
coefficients are positive; all eventual divisions by an index are valid because
both indices converge to \(\alpha_0>0\).

Let
\[
q_{1,n}=F_{1,n}^{\leftarrow}(1-1/t_n).
\tag{7}
\]
Continuity in \(\mathrm{ass{:}continuous\text{-}margins}\) gives
\[
P_{\theta,n}(X_{1,n}>q_{1,n})=t_n^{-1},
\tag{8}
\]
and it makes the event in (8) equal, up to a null set, to
\(\{F_{1,n}(X_{1,n})>1-1/t_n\}\). By
\(\mathrm{ass{:}local\text{-}tail\text{-}collision}\), there are constants
\(0<\underline\alpha<\overline\alpha<\infty\) such that, eventually,
\[
\underline\alpha\leq\alpha_{j,n}\leq\overline\alpha,
\qquad j\in\{1,2\}.
\tag{9}
\]
Fix \(d\in(0,1)\). If \(M_n=a_n\), then
\(\mathrm{ass{:}uniform\text{-}tail\text{-}normalization}\), evaluated at the
fixed multiplier \(d\), gives
\[
t_nP_{\theta,n}(\epsilon_{1,n}>dM_n)
=d^{-\alpha_{1,n}}\{1+o(1)\}>1
\tag{10}
\]
eventually. If \(M_n=b_n\), the same conclusion holds for
\(P_{\theta,n}(\beta_n\epsilon_{2,n}>dM_n)\). Positivity in (1) and (8)
therefore force \(q_{1,n}\geq dM_n\).

Choose a fixed \(C>2\) so large that
\(2(C/2)^{-\underline\alpha}<1\). By (6), each of
\(CM_n/(2a_n)\) and \(CM_n/(2b_n)\) lies in a fixed compact subset of
\((0,\infty)\), and each is at least \(C/2\). The union bound and
\(\mathrm{ass{:}uniform\text{-}tail\text{-}normalization}\) give
\[
\begin{aligned}
P_{\theta,n}(X_{1,n}>CM_n)
&\leq P_{\theta,n}(\epsilon_{1,n}>CM_n/2)
+P_{\theta,n}(\beta_n\epsilon_{2,n}>CM_n/2),\\
t_nP_{\theta,n}(X_{1,n}>CM_n)
&\leq2(C/2)^{-\underline\alpha}\{1+o(1)\}<1
\end{aligned}
\tag{11}
\]
eventually. Equation (8) then forces \(q_{1,n}\leq CM_n\). Thus
\[
dM_n\leq q_{1,n}\leq CM_n.
\tag{12}
\]
In particular, the two ratios \(q_{1,n}/a_n\) and \(q_{1,n}/b_n\) are
bounded above and bounded away from zero.

Define the propagated and idiosyncratic component exceedance probabilities
\[
p_{P,n}=P_{\theta,n}(\beta_n\epsilon_{2,n}>q_{1,n}),
\qquad
p_{I,n}=P_{\theta,n}(\epsilon_{1,n}>q_{1,n}).
\tag{13}
\]
Equations (3), (6), (12), and
\(\mathrm{ass{:}uniform\text{-}tail\text{-}normalization}\), applied at the
two compact multiplier sequences, give
\[
\begin{aligned}
p_{P,n}
&=t_n^{-1}\left(\frac{q_{1,n}}{b_n}\right)^{-\alpha_{2,n}}
\{1+o(1)\},\\
p_{I,n}
&=t_n^{-1}\left(\frac{q_{1,n}}{a_n}\right)^{-\alpha_{1,n}}
\{1+o(1)\}.
\end{aligned}
\tag{14}
\]
Hence \(p_{P,n}\asymp t_n^{-1}\) and \(p_{I,n}\asymp t_n^{-1}\).

We next derive, rather than import through a relabeled model class, the required
one-large-jump expansion. Put \(Y_n=\beta_n\epsilon_{2,n}\) and
\(W_n=\epsilon_{1,n}\). For fixed \(\delta\in(0,1/2)\), positivity implies
\[
\begin{aligned}
p_{P,n}+p_{I,n}-p_{P,n}p_{I,n}
&\leq P_{\theta,n}(Y_n+W_n>q_{1,n})\\
&\leq P_{\theta,n}(Y_n>(1-\delta)q_{1,n})
+P_{\theta,n}(W_n>(1-\delta)q_{1,n})\\
&\quad+P_{\theta,n}(Y_n>\delta q_{1,n})
P_{\theta,n}(W_n>\delta q_{1,n}).
\end{aligned}
\tag{15}
\]
The product identities in (15) use
\(\mathrm{ass{:}innovation\text{-}independence}\). By (12), every threshold
in (15), normalized by the relevant tail quantile, lies in a fixed compact
subset of \((0,\infty)\). Thus tail normalization and (14) give
\[
\frac{P_{\theta,n}(Y_n>(1-\delta)q_{1,n})}{p_{P,n}}
=(1-\delta)^{-\alpha_{2,n}}+o(1),
\tag{16}
\]
with the analogous identity for \(W_n\), while the final product in (15) is
\(O(t_n^{-2})\). Assumption \(\mathrm{ass{:}intermediate\text{-}sequence}\)
gives \(t_n=n/k_n\to\infty\); therefore both products in (15) are
\(o(p_{P,n}+p_{I,n})\). Divide (15) by \(p_{P,n}+p_{I,n}\), use (9), and
then let \(\delta\downarrow0\). The result is
\[
P_{\theta,n}(X_{1,n}>q_{1,n})
=p_{P,n}+p_{I,n}+o(p_{P,n}+p_{I,n}).
\tag{17}
\]

Set \(H_n=F_{2,n}(X_{2,n})=F_{2,n}(\epsilon_{2,n})\). To verify its law
directly, for \(z\in(0,1)\) take
\(x_z=F_{2,n}^{\leftarrow}(z)\). Continuity gives \(F_{2,n}(x_z)=z\), and
the events \(\{F_{2,n}(X_{2,n})\leq z\}\) and \(\{X_{2,n}\leq x_z\}\)
differ only on a null interval on which the distribution assigns no mass.
Consequently \(P_{\theta,n}(H_n\leq z)=z\), so
\[
0\leq H_n\leq1,
\qquad
E_{P_{\theta,n}}(H_n)=\frac12.
\tag{18}
\]
On \(\{Y_n>q_{1,n}\}\), monotonicity of \(F_{2,n}\) gives
\[
0\leq1-H_n
\leq P_{\theta,n}(\epsilon_{2,n}>q_{1,n}/\beta_n)=p_{P,n}.
\tag{19}
\]
Therefore
\[
E[H_n\mathbf 1\{Y_n>q_{1,n}\}]
=p_{P,n}+O(p_{P,n}^2),
\qquad
E[H_n\mathbf 1\{W_n>q_{1,n}\}]
=\frac12p_{I,n},
\tag{20}
\]
where the second equality uses innovation independence and (18). The
intersection of the two component exceedance events has probability
\(p_{P,n}p_{I,n}=o(p_{P,n}+p_{I,n})\). By (17), the part of
\(\{Y_n+W_n>q_{1,n}\}\) outside their union also has probability
\(o(p_{P,n}+p_{I,n})\). Since \(0\leq H_n\leq1\), equations (17)--(20)
and the conditioning-event equality after (8) prove
\[
\Gamma_{12,n}(t_n)
=\frac{p_{P,n}+p_{I,n}/2}{p_{P,n}+p_{I,n}}+o(1)
=\frac12+\frac12\frac{p_{P,n}/p_{I,n}}{1+p_{P,n}/p_{I,n}}+o(1).
\tag{21}
\]

The ratio in (14) satisfies
\[
\frac{p_{P,n}}{p_{I,n}}
=\beta_n^{\alpha_{2,n}}\frac{c_{2,n}}{c_{1,n}}
q_{1,n}^{\Delta_n}\{1+o(1)\}.
\tag{22}
\]
The intrinsic-coordinate hypothesis, (5), the coefficient limit, and
\(\Delta_n\to0\) give
\[
\Delta_n\log a_n=-\xi_{\mathrm{rev},n}\to-\xi_{\mathrm{rev}},
\qquad
\Delta_n\log b_n
=\Delta_n\log a_n
+\Delta_n\log\frac{u_{2,n}}{u_{1,n}}
+\Delta_n\log\beta_n
\to-\xi_{\mathrm{rev}}.
\tag{23}
\]
It follows that \(\Delta_n\log M_n\to-\xi_{\mathrm{rev}}\), irrespective
of which term realizes the maximum. Equation (12) then gives
\[
\Delta_n\log q_{1,n}\longrightarrow-\xi_{\mathrm{rev}},
\tag{24}
\]
because \(\log(q_{1,n}/M_n)\) is bounded and \(\Delta_n\to0\). Finally,
\(\beta_n^{\alpha_{2,n}}\to\beta_0^{\alpha_0}\) and
\(c_{2,n}/c_{1,n}\to c_0^{-1}\), so (22)--(24) imply
\[
\frac{p_{P,n}}{p_{I,n}}
\longrightarrow
\beta_0^{\alpha_0}c_0^{-1}e^{-\xi_{\mathrm{rev}}}
=R_{\mathrm{rev},\xi}.
\tag{25}
\]
Equations (21) and (25) prove the asserted limit for
\(\Gamma_{12,n}(t_n)\).

For the other rank coefficient, reverse orientation gives
\(F_{2,n}^{\leftarrow}(1-1/t_n)=u_{2,n}\). On
\(\{X_{2,n}>u_{2,n}\}\), positivity in (1) gives
\(X_{1,n}>\beta_n u_{2,n}=b_n\). Hence monotonicity and continuity imply
\[
\begin{aligned}
0\leq1-\Gamma_{21,n}(t_n)
&\leq1-F_{1,n}(b_n)
=P_{\theta,n}(X_{1,n}>b_n)\\
&\leq P_{\theta,n}(\epsilon_{2,n}>u_{2,n}/2)
+P_{\theta,n}(\epsilon_{1,n}>b_n/2).
\end{aligned}
\tag{26}
\]
The first probability in the last line is \(O(t_n^{-1})\). By (6), the
second threshold is a compact positive multiple of \(u_{1,n}\), so the
second probability is also \(O(t_n^{-1})\), by uniform tail normalization.
Since \(t_n\to\infty\), (26) proves
\[
\Gamma_{21,n}(t_n)\longrightarrow1.
\tag{27}
\]
Combining (21), (25), and (27) with
\(D_n(t_n)=\Gamma_{12,n}(t_n)-\Gamma_{21,n}(t_n)\) from
\(\mathrm{def{:}finite\text{-}threshold\text{-}ctc}\) gives
\[
D_n(t_n)\longrightarrow-\frac{1}{2(1+R_{\mathrm{rev},\xi})}.
\tag{28}
\]
This entire derivation used the original reverse array and the original
balance limit \(c_0\); in particular, it did not assert membership of a
label-exchanged array in \(\Theta_{\mathrm{CTC}}(B)\).

For the finite logarithmic-coordinate corollary, assume additionally
\(\Delta_n\log c_{2,n}\to0\) and \(\lambda_n\to\lambda\in\mathbb R\).
Tail-constant balance makes \(\log(c_{1,n}/c_{2,n})\) bounded, and local
collision gives \(\Delta_n\to0\). Therefore
\[
\Delta_n\log c_{1,n}
=\Delta_n\log c_{2,n}
+\Delta_n\log(c_{1,n}/c_{2,n})
\longrightarrow0.
\tag{29}
\]
The \(j=1\) identity in (3) now gives the exact equality
\[
\xi_{\mathrm{rev},n}
=-\frac{\lambda_n+\Delta_n\log c_{1,n}}{\alpha_{1,n}}
\longrightarrow-\frac{\lambda}{\alpha_0}.
\tag{30}
\]
Substituting (30) into (25) yields
\[
R_{\mathrm{rev},\xi}
=\beta_0^{\alpha_0}c_0^{-1}e^{\lambda/\alpha_0}
=R_{\mathrm{rev}}(\lambda).
\tag{31}
\]
By the definition of \(g_{\alpha_0,\beta_0,1/c_0}\),
\[
-\frac{1}{2\{1+R_{\mathrm{rev}}(\lambda)\}}
=-g_{\alpha_0,\beta_0,1/c_0}(-\lambda),
\tag{32}
\]
which proves every finite-\(\lambda\) corollary in the statement.

It remains to prove the two extended exact-Pareto endpoints without importing
a forward theorem under a relabeled balance constant. On the stated reverse
scaled exact-Pareto subclass, put
\[
s_{j,n}=c_{j,n}^{1/\alpha_{j,n}},
\qquad
Y_n=\beta_n\epsilon_{2,n},
\qquad
W_n=\epsilon_{1,n}.
\tag{33}
\]
Thus \(s_{j,n}\) is the lower support endpoint and
\[
P(Y_n>v)=c_{2,n}\beta_n^{\alpha_{2,n}}v^{-\alpha_{2,n}},
\qquad
P(W_n>v)=c_{1,n}v^{-\alpha_{1,n}}
\tag{34}
\]
whenever \(v\) exceeds the corresponding propagated or idiosyncratic support
endpoint. Define
\[
a_n^{\mathrm{Par}}=\beta_n s_{2,n}t_n^{1/\alpha_{2,n}},
\qquad
b_n^{\mathrm{Par}}=s_{1,n}t_n^{1/\alpha_{1,n}},
\qquad
M_n^{\mathrm{Par}}=\max(a_n^{\mathrm{Par}},b_n^{\mathrm{Par}}).
\tag{35}
\]
The variables in (35) are the component \((1-1/t_n)\)-quantiles. If
\(q_{1,n}^{\mathrm{Par}}\) is the \((1-1/t_n)\)-quantile of \(Y_n+W_n\),
positivity gives \(q_{1,n}^{\mathrm{Par}}\geq M_n^{\mathrm{Par}}\). Conversely,
the union bound at the two component \((1-1/(2t_n))\)-quantiles gives
\[
q_{1,n}^{\mathrm{Par}}
\leq\beta_n s_{2,n}(2t_n)^{1/\alpha_{2,n}}
+s_{1,n}(2t_n)^{1/\alpha_{1,n}}
\leq C_{\mathrm{Par}}M_n^{\mathrm{Par}}
\tag{36}
\]
for a fixed finite \(C_{\mathrm{Par}}\), using the index bounds in (9).
Moreover,
\[
\frac{q_{1,n}^{\mathrm{Par}}}{\beta_n s_{2,n}+s_{1,n}}
\geq\frac12\min_{j\in\{1,2\}}t_n^{1/\alpha_{j,n}}
\longrightarrow\infty.
\tag{37}
\]
Thus both component exceedance probabilities at
\(q_{1,n}^{\mathrm{Par}}\) tend to zero. Write them as
\(p_{P,n}^{\mathrm{Par}}\) and \(p_{I,n}^{\mathrm{Par}}\). For fixed
\(\delta\in(0,1/2)\), the exact tails give
\[
P(Y_n>(1-\delta)q_{1,n}^{\mathrm{Par}})
=(1-\delta)^{-\alpha_{2,n}}p_{P,n}^{\mathrm{Par}},
\qquad
P(W_n>(1-\delta)q_{1,n}^{\mathrm{Par}})
=(1-\delta)^{-\alpha_{1,n}}p_{I,n}^{\mathrm{Par}},
\tag{38a}
\]
and analogous equalities with \(\delta\) in place of \(1-\delta\).
Moreover,
\[
0\leq
\frac{p_{P,n}^{\mathrm{Par}}p_{I,n}^{\mathrm{Par}}}
{p_{P,n}^{\mathrm{Par}}+p_{I,n}^{\mathrm{Par}}}
\leq\min(p_{P,n}^{\mathrm{Par}},p_{I,n}^{\mathrm{Par}})\longrightarrow0.
\tag{38b}
\]
Insert (38a)--(38b) into the deterministic event bounds in (15), divide by
\(p_{P,n}^{\mathrm{Par}}+p_{I,n}^{\mathrm{Par}}\), and then let
\(\delta\downarrow0\), using \(\alpha_{j,n}\to\alpha_0\). This proves,
uniformly over the component-probability ratio,
\[
P(Y_n+W_n>q_{1,n}^{\mathrm{Par}})
=p_{P,n}^{\mathrm{Par}}+p_{I,n}^{\mathrm{Par}}
+o(p_{P,n}^{\mathrm{Par}}+p_{I,n}^{\mathrm{Par}}),
\tag{38}
\]
where
\[
\frac{p_{P,n}^{\mathrm{Par}}}{p_{I,n}^{\mathrm{Par}}}
=\beta_n^{\alpha_{2,n}}\frac{c_{2,n}}{c_{1,n}}
(q_{1,n}^{\mathrm{Par}})^{\Delta_n}.
\tag{39}
\]
The rank-numerator calculation (18)--(21), with exact-Pareto tails in place of
tail normalization, therefore yields
\[
\Gamma_{12,n}(t_n)
=\frac12+\frac12
\frac{p_{P,n}^{\mathrm{Par}}/p_{I,n}^{\mathrm{Par}}}
{1+p_{P,n}^{\mathrm{Par}}/p_{I,n}^{\mathrm{Par}}}+o(1).
\tag{40}
\]

Under \(\Delta_n\log c_{2,n}\to0\), equation (29) remains valid. Equations
(35) then imply
\[
\begin{aligned}
\Delta_n\log a_n^{\mathrm{Par}}
&=\frac{\lambda_n}{\alpha_{2,n}}+o(1),\\
\Delta_n\log b_n^{\mathrm{Par}}
&=\frac{\lambda_n}{\alpha_{1,n}}+o(1).
\end{aligned}
\tag{41}
\]
If \(\lambda_n\to\lambda\in[-\infty,+\infty]\), both right sides have
limit \(\lambda/\alpha_0\) for finite \(\lambda\) and the corresponding
signed infinite limit at the endpoints. Equations (36),
\(\Delta_n\to0\), and the elementary maximum case split therefore give
\[
\Delta_n\log q_{1,n}^{\mathrm{Par}}
\longrightarrow
\begin{cases}
-\infty,&\lambda=-\infty,\\
\lambda/\alpha_0,&\lambda\in\mathbb R,\\
+\infty,&\lambda=+\infty.
\end{cases}
\tag{42}
\]
Substitution into (39), together with the coefficient and tail-balance limits,
gives
\[
\frac{p_{P,n}^{\mathrm{Par}}}{p_{I,n}^{\mathrm{Par}}}
\longrightarrow
\begin{cases}
0,&\lambda=-\infty,\\
R_{\mathrm{rev}}(\lambda),&\lambda\in\mathbb R,\\
+\infty,&\lambda=+\infty.
\end{cases}
\tag{43}
\]

Finally, the exact second-margin threshold is
\(U_{2,n}(t_n)=s_{2,n}t_n^{1/\alpha_{2,n}}\). Direct algebra gives
\[
\log\frac{U_{2,n}(t_n)}{s_{1,n}}
=\frac{\log t_n}{\alpha_{2,n}}
+\frac{\Delta_n\log c_{2,n}}{\alpha_{1,n}\alpha_{2,n}}
-\frac{\log(c_{1,n}/c_{2,n})}{\alpha_{1,n}}
\longrightarrow+\infty.
\tag{44}
\]
On \(\{X_{2,n}>U_{2,n}(t_n)\}\), one has
\(X_{1,n}>\beta_nU_{2,n}(t_n)\). Hence, exactly as in (26),
\[
\begin{aligned}
0\leq1-\Gamma_{21,n}(t_n)
&\leq P(X_{1,n}>\beta_nU_{2,n}(t_n))\\
&\leq P(\epsilon_{2,n}>U_{2,n}(t_n)/2)
+P(\epsilon_{1,n}>\beta_nU_{2,n}(t_n)/2)\longrightarrow0.
\end{aligned}
\tag{45}
\]
Indeed, the first probability is \(2^{\alpha_{2,n}}/t_n\to0\), and the
second tends to zero by (44), the coefficient limit, the positive lower index
bound in (9), and the exact tail formula. Equations (40), (43), and (45),
followed by subtraction in the definition of \(D_n(t_n)\), give respectively
\[
(\Gamma_{21,n},\Gamma_{12,n},D_n)
\longrightarrow(1,\tfrac12,-\tfrac12),
\quad
\left(1,\frac12+\frac{R_{\mathrm{rev}}(\lambda)}
{2\{1+R_{\mathrm{rev}}(\lambda)\}},
-g_{\alpha_0,\beta_0,1/c_0}(-\lambda)\right),
\quad
(1,1,0)
\tag{46}
\]
in the cases \(\lambda=-\infty\), finite \(\lambda\), and
\(\lambda=+\infty\), respectively. Here every displayed rank coefficient
in (46) is evaluated at \(t_n\). This proves the extended exact-Pareto rows and
completes the proof.
\end{proof}
\par\smallskip\noindent \textit{Justification.} A direct calculation on the original reverse array controls the marginal threshold by the two component quantiles, proves a ratio-free one-large-jump expansion, and evaluates the conditional ranks under the original balance limit \(c_0\); direct exact-Pareto algebra then gives the logarithmic-coordinate formula and both endpoints. \textit{Gap.} The reverse class curve now covers every finite intrinsic threshold collision; no class-uniform remainder or inference guarantee is claimed. \textit{Consumer.} It supplies the reverse-orientation population sign and magnitude from the original \(c_0\)-balanced array, without treating exact label exchange as a proof mechanism, exchanged-class membership, or observed-law orientation identification.

\begin{theorem}[thm:four-coordinate-ambiguity-obstruction]\label{thm:four-coordinate-ambiguity-obstruction}
Fix \(\kappa\in[0,\infty]\) and a compact \(K\subset\mathcal L_\kappa\) containing \((0,0,0,0)\).  Under the frozen collision model, the proposed local experiment indexed only by \((\lambda,h,b_1,b_2)\), together with the summed-marginal information form
\[
 I_{rs}=\sum_{a=1}^2\int\dot\ell_{a,r}(z)\dot\ell_{a,s}(z)\,dM_{\vartheta_0}^{(a)}(z),
\]
is invalid for two distinct reasons.  First, the forward zero-bias exact-Pareto subclass contains two arrays with
\[
 (\lambda_n,h_n,b_{1,n},b_{2,n})=(0,0,0,0)
\]
whose restrictions to \(\mathcal T_n\) have total-variation distance tending to one.  Hence those four coordinates do not index a single local law; the omitted distinguishing coordinate is a root-\(k_n\) common-tail-index coordinate.  Second, even at a fixed common index, \(N_n^{(1)}\) and \(N_n^{(2)}\) have order-\(k_n\) joint exceedance overlap.  After root-\(k_n\) normalization this overlap contributes a nonzero cross-covariance term to the joint score variance, a term absent from \(I_{rs}\).  Thus \(I_{rs}\) is not, in general, the information form of the joint tail-rank experiment.  Any valid replacement must at least add the common-index coordinate and a joint overlap intensity, or an equivalent paired marked point process.  This obstruction does not construct the augmented experiment, a nonconstant honest direction set, or a matching ambiguity frontier.
\end{theorem}
\begin{proof}
Let \(v_0=(0,0,0,0)\).  Because \(v_0\in K\), the first certificate in \(\mathrm{thm{:}joint\text{-}collision\text{-}experiment\text{-}obstruction}\) supplies two forward zero-bias exact-Pareto arrays, indexed below by \(r\in\{0,1\}\), such that
\[
 (\lambda_n^{(r)},h_n^{(r)},b_{1,n}^{(r)},b_{2,n}^{(r)})=v_0
 \qquad(r=0,1)
 \tag{1}
\]
for every \(n\), while, with
\[
 P_{r,n}^{\mathcal T}
 :=(P_{\theta^{(r)},n}^{\otimes n})\big|_{\mathcal T_n},
\]
their tail-rank restrictions satisfy
\[
 \left\|P_{0,n}^{\mathcal T}-P_{1,n}^{\mathcal T}\right\|_{\mathrm{TV}}
 \longrightarrow1.
 \tag{2}
\]
All measures in (2) are probability measures on the common measurable space \(\mathcal T_n\), by \(\mathrm{def{:}tail\text{-}rank\text{-}experiment}\), so the total-variation distance is well typed.

To make the indexing contradiction explicit, suppose that a local experiment indexed only by \(v=(\lambda,h,b_1,b_2)\) represented both arrays uniformly, and denote its law at \(v_0\) by \(Q_{n,v_0}\).  Representation at their common coordinate (1) would require
\[
 \left\|P_{r,n}^{\mathcal T}-Q_{n,v_0}\right\|_{\mathrm{TV}}
 \longrightarrow0
 \qquad(r=0,1).
 \tag{3}
\]
The triangle inequality for total variation then gives
\[
 \left\|P_{0,n}^{\mathcal T}-P_{1,n}^{\mathcal T}\right\|_{\mathrm{TV}}
 \leq
 \left\|P_{0,n}^{\mathcal T}-Q_{n,v_0}\right\|_{\mathrm{TV}}
 +
 \left\|Q_{n,v_0}-P_{1,n}^{\mathcal T}\right\|_{\mathrm{TV}}
 \longrightarrow0,
 \tag{4}
\]
contradicting (2).  Thus the four coordinates do not index a single local law.  The Hill certificate in \(\mathrm{thm{:}joint\text{-}collision\text{-}experiment\text{-}obstruction}\) identifies the variation suppressed by (1) as root-\(k_n\) movement of the common tail index.  Therefore a common-tail-index coordinate is necessary, although the cited obstruction theorem does not assert that adding it is sufficient.

We next isolate the independent information-form failure.  Whenever the marginal intensity scores in \(\mathrm{def{:}joint\text{-}experiment\text{-}handle}\) exist, define their centered contributions by
\[
 S_{a,n}(r)
 :=k_n^{-1/2}\int\dot\ell_{a,r}(z)
 \{dN_n^{(a)}(z)-k_n\,dM_{\vartheta_0}^{(a)}(z)\},
 \qquad a\in\{1,2\}.
 \tag{5}
\]
The positivity of \(k_n\) makes the normalization in (5) valid.  For every direction for which the second moments in (5) are finite, the elementary variance identity is
\[
 \operatorname{Var}\{S_{1,n}(r)+S_{2,n}(r)\}
 =\operatorname{Var}\{S_{1,n}(r)\}
 +\operatorname{Var}\{S_{2,n}(r)\}
 +2\operatorname{Cov}\{S_{1,n}(r),S_{2,n}(r)\}.
 \tag{6}
\]
The second certificate in \(\mathrm{thm{:}joint\text{-}collision\text{-}experiment\text{-}obstruction}\) states that, even at a fixed common index, the two point measures share joint exceedances with intensity of order \(k_n\), and that the corresponding normalized overlap covariance is nonzero.  Since each contribution in (5) is scaled by \(k_n^{-1/2}\), an order-\(k_n\) covariance before normalization is an order-one cross term in (6).  In contrast, the displayed candidate
\[
 I_{rs}=\sum_{a=1}^2\int\dot\ell_{a,r}(z)\dot\ell_{a,s}(z)\,dM_{\vartheta_0}^{(a)}(z)
 \tag{7}
\]
contains only the two marginal terms and no cross term.  Equations (6)--(7) and the nonzero-overlap certificate therefore show that (7) cannot in general equal the joint score covariance, and hence cannot be the information form of the joint tail-rank experiment.

The contradictions in (4) and (6)--(7) are logically separate: adding a common-index coordinate does not restore the missing overlap term, and recording a joint overlap intensity does not make the four coordinates distinguish the arrays in (1).  Consequently a valid replacement requires at least both augmentations stated in the theorem.  No likelihood expansion, quadratic-mean differentiability, confidence-set calibration, or opposite-orientation contiguity argument is supplied by either certificate, so the augmented experiment, a nonconstant honest set, and a matching ambiguity frontier remain outside the conclusion.
\end{proof}
\par\smallskip\noindent \textit{Justification.} The theorem certifies two independent failures of the proposed local handle: four coordinates suppress root-\(k_n\) common-index variation, and the summed marginal information suppresses the order-\(k_n\) joint exceedance overlap. \textit{Gap.} The frozen primitives do not construct the augmented paired marked process, a uniform likelihood or quadratic-mean expansion, a calibrated nonconstant honest set, or an opposite-orientation contiguity family. \textit{Consumer.} The result prevents the four-coordinate summed-marginal handle from being used as a valid joint experiment or as the basis of a claimed ambiguity frontier, and identifies two necessary augmentations.

\section{Interpretation}

The population collision curves describe how a chosen intermediate threshold changes the finite-threshold CTC contrast; they do not identify structural orientation from the unrestricted observational law.  The experiment obstruction is specific: it shows that the proposed four coordinates and a sum of marginal intensity informations are insufficient.  It does not prove that inference is impossible after adding a common-index coordinate and a paired overlap process.

\section*{Technical internal limitation (diagnostic only)}

The structural orientation is a component of \(\theta_n\), not an asserted functional of \(P_{\theta,n}\), and no unique structural representation is assumed.  The current inference claims are restricted to \(\mathcal T_n\).  The frozen primitives provide neither an augmented joint likelihood or quadratic-mean expansion, nor a contrast-bias calibration, nor a bounded-divergence opposite-orientation construction.  Accordingly, no positive ambiguity lower bound, matching frontier, or vanishing-ambiguity result is asserted for either the tail-rank experiment or unrestricted full-distribution causal discovery.

\section{Honest scope}

The proved population results are the scaled exact-Pareto and class-level forward and reverse finite-threshold collision curves, the structural witness membership result, and the exact scale geometry.  The proved experiment result is negative: the four-coordinate handle omits root-\(k_n\) common-index variation, and the summed marginal information omits the nonzero order-\(k_n\) overlap covariance.  The only certified confidence-set construction is the constant full-alphabet family, which has exact coverage, label equivariance, envelope-grid nesting, and ambiguity probability one.  Honesty is conditional on the declared envelope \(B_n\); the assumptions do not identify or estimate that envelope.  Tangent existence, paired conditional rank-mark kernels, the full bias drift, a uniform likelihood remainder, an augmented paired experiment, a nonconstant envelope-conditional direction set, and an opposite-orientation contiguity argument establishing a matching ambiguity frontier remain open.

\begin{thebibliography}{99}

  \bibitem{LeimenstollSchienle2026} Lisa Leimenstoll and Melanie Schienle. Identification and Inference for Causal Effects in Extremes under General Conditions. arXiv:2608.22957v1, 2026. Proposition 3.1; Theorems 4.1, 4.3, and 4.4; Appendix A.6; Sections 5--7.

  \bibitem{GneccoEtAl2021} Nicola Gnecco, Nicolai Meinshausen, Jonas Peters, and Sebastian Engelke. Causal discovery in heavy-tailed models. Annals of Statistics, 49(3):1755--1778, 2021. Lemma 1 and Theorems 1--2.

  \bibitem{PascheChavezDemoulinDavison2023} Olivier C. Pasche, Valérie Chavez-Demoulin, and Anthony C. Davison. Causal modelling of heavy-tailed variables and confounders with application to river flow. Extremes, 26:573--594, 2023. Sections 2--3 and 5; doi:10.1007/s10687-022-00456-4.

  \bibitem{BodikPalusPawlas2024} Juraj Bodik, Milan Paluš, and Zbyněk Pawlas. Causality in extremes of time series. Extremes, 27:67--121, 2024. Definition 1, Theorems 1--2, Section 4, and Section 6; doi:10.1007/s10687-023-00479-5.

  \bibitem{BeirlantBouquiauxWerker2006} Jan Beirlant, Christel Bouquiaux, and Bas J. M. Werker. Semiparametric lower bounds for tail-index estimation. Journal of Statistical Planning and Inference, 136(3):705--729, 2006. Theorems 2.1--2.2.

  \bibitem{Drees1998} Holger Drees. On smooth statistical tail functionals. Scandinavian Journal of Statistics, 25(1):187--210, 1998. Theorem 2.1.

  \bibitem{BuecherDette2013} Axel B\"ucher and Holger Dette. Multiplier bootstrap of tail copulas with applications. Bernoulli, 19(5A):1655--1687, 2013. Theorem 2.2 and Theorems 3.2--3.4.

  \bibitem{CarpentierKim2014CI} Alexandra Carpentier and Arlene K. H. Kim. Adaptive confidence intervals for the tail coefficient in a wide second order class of Pareto models. Electronic Journal of Statistics, 8(2):2067--2110, 2014. Theorems 3.2 and 3.6.

  \bibitem{SasakiWang2024} Yuya Sasaki and Yulong Wang. Uniform confidence intervals for the tail index and the extreme quantile. Journal of Econometrics, 241(1):105865, 2024. Impossibility and uniform-inference theorems.

  \bibitem{Novak2014} S. Y. Novak. Lower bounds to the accuracy of inference on heavy tails. Bernoulli, 20(2):979--989, 2014. Main minimax lower-bound theorems.

  \bibitem{deHaanFerreira2006} Laurens de Haan and Ana Ferreira. Extreme Value Theory: An Introduction. Springer, 2006. Sections 2.3--2.4.

  \bibitem{vanDerVaart1998} A. W. van der Vaart. Asymptotic Statistics. Cambridge University Press, 1998. Chapters 6--9.

\end{thebibliography}

\end{document}